\documentclass[11pt,a4paper]{article}

%─── Packages ─────────────────────────────────────────────────────────────────
\usepackage{amsmath,amssymb,amsthm}
\usepackage{geometry}
\usepackage{tabularray}
\UseTblrLibrary{booktabs}
\usepackage{booktabs}
\usepackage{array}
\usepackage{microtype}
\usepackage{hyperref}
\usepackage{xcolor}
\usepackage{enumitem}
\geometry{top=2.5cm,bottom=2.8cm,left=2.5cm,right=2.5cm}

%─── Theorem environments ─────────────────────────────────────────────────────
\newtheorem{theorem}{Theorem}[section]
\newtheorem{lemma}[theorem]{Lemma}
\newtheorem{proposition}[theorem]{Proposition}
\newtheorem{corollary}[theorem]{Corollary}
\theoremstyle{definition}
\newtheorem{definition}[theorem]{Definition}
\newtheorem{result}[theorem]{Result}
\theoremstyle{remark}
\newtheorem{remark}[theorem]{Remark}

%─── Macros ───────────────────────────────────────────────────────────────────
\newcommand{\vp}{\varphi}
\providecommand{\ms}{\mu^{*}}
\newcommand{\bst}{\beta^{*}}
\providecommand{\Kfb}{K_{\mathrm{fb}}}
\providecommand{\Efb}{E_{\mathrm{fb}}}
\newcommand{\lam}{\lambda}
\newcommand{\bz}{\beta_{0}}
\newcommand{\TCMB}{T_{\mathrm{CMB}}}
\newcommand{\rCMB}{\rho_{\mathrm{CMB}}}
\newcommand{\Lcond}{\Lambda_{\mathrm{cond}}}
\newcommand{\vEW}{v_{\mathrm{EW}}}
\newcommand{\mWZW}{m_{\mathrm{WZW}}}
\newcommand{\aInv}{\alpha^{-1}}
\newcommand{\sinW}{\sin^{2}\!\theta_{W}}
\newcommand{\MPl}{M_{\mathrm{Pl}}}
\newcommand{\LUV}{\Lambda_{\mathrm{UV}}}
\newcommand{\Lobs}{\Lambda_{\mathrm{obs}}}
\newcommand{\rinfty}{\rho_{\infty}}
\newcommand{\oBD}{\omega_{\mathrm{BD}}}
\newcommand{\xNMC}{\xi_{\mathrm{NMC}}}
\newcommand{\Ja}{J_{2a}}
\newcommand{\Jfb}{J_{2\mathrm{iso}}^{\mathrm{fb}}}
\newcommand{\Jfour}{J_{4}}

%─── Colours ──────────────────────────────────────────────────────────────────
\definecolor{darkgreen}{RGB}{0,120,60}
\definecolor{darkred}{RGB}{160,20,20}
\definecolor{darkorange}{RGB}{180,80,0}

%══════════════════════════════════════════════════════════════════════════════
\title{\textbf{Emergent Gravity, Inflation, Dark Energy,\\
and the Weak Equivalence Principle\\
from the Density-Feedback Hopfion Condensate}
\\[0.6em]
{\large Companion Paper VII to the Density-Feedback
Faddeev--Niemi Hopfion Series}
\\[0.6em]
\large Preprint --- April 2026}
\author{Frederick Manfredi}
\date{Superseded estimate update - June 2026}
%══════════════════════════════════════════════════════════════════════════════

\begin{document}
\maketitle

\paragraph*{Units and conventions.}
Natural units $\hbar = c = k_B = 1$ throughout.
Energies and masses are in eV, MeV, or GeV as appropriate.
The Planck mass is the \emph{unreduced} value
$\MPl = 1.22\times10^{19}\,\mathrm{GeV}$ consistently throughout
Papers~I--VII (the reduced value $\MPl^{\mathrm{red}} = \MPl/\sqrt{8\pi}
\approx 2.435\times10^{18}\,\mathrm{GeV}$ is not used).
The condensate scale is $\Lcond = \TCMB(\pi^2/150)^{1/4} = 1.1895\times10^{-4}\,\mathrm{eV}$.

\begin{abstract}
We extend the density-feedback Faddeev--Niemi Hopfion
framework of~\cite{Paper1,Paper2,Paper3,Paper4,Paper5,Paper6} from its
Standard-Model and QCD predictions to cosmological and gravitational
observables.  Starting from four axioms of the semi-Dirac scalar
field $\rho$ whose condensate is the $Q=2$ icosahedral Hopfion
vacuum established in~\cite{Paper1}, we derive:
Newtonian gravity from gradient flux imbalance with $G_N =
G_{\mathrm{src}}^2/(4\pi\vp^6)$; the weak-field Schwarzschild and
Kerr acoustic metrics with $G_{\mu\nu}=0$ in the exterior verified
without assumption; an induced Einstein--Hilbert term from one-loop
fluctuations with $\MPl^2 = \LUV^2/(32\pi^2)$ and
$\LUV\approx\MPl/\vp^6\approx0.0557\,\MPl$ (exact: $\LUV=\MPl/(4\pi\sqrt{2})$,
within 1\% of $\MPl/\vp^6$ by the Bogomolny coincidence $4\pi\sqrt{2}\approx\vp^6$);
$\Lobs = \Lambda_0 e^{-\vp^7/\beta}\approx10^{-122}\MPl^4$; CMB
inflationary predictions $n_s = 0.9654$ (within $0.12\sigma$ of
Planck~2018) and $r=0.0036$, with exact normalisation
$\mathcal{P}_s=2.10\times10^{-9}$ from gravitational reheating at $N_e=57.73$;
and the exact Hawking temperature with a soup correction
$\delta T_H/T_H\approx+4.2\times10^{-5}$.
The full dynamical Einstein equations are derived from the
condensate effective action, with the frozen-attractor limit
yielding $G_{\mu\nu}+\Lambda_{\rm phys}g_{\mu\nu}=0$ exactly,
where $F(\rinfty)=\MPl^2(39\vp+25)/2$ is an exact golden-ratio
consequence; energy-momentum conservation $\nabla^\mu T_{\mu\nu}=0$
follows from scalar field equation consistency (Bianchi identity).
The exact rotating condensate background is given by the Kerr mass
multipole series $\Phi_{\rm Kerr}=M/r-Ma^2P_2/(2r^3)+\ldots$,
satisfying $\nabla^2\Phi_{\rm Kerr}=0$ exactly term by term.
The Cassini PPN constraint is proved resolved by Vainshtein screening,
with $r_V\approx9.3\times10^6\,\mathrm{AU}$ derived without free parameters.
The condensate scalar field acts as dark matter with exact golden-ratio
sound speed $c_s=1/\vp$ (Theorem~\ref{P7:thm:cs}), is pressureless on all
galactic scales ($\lambda_J\approx2\times10^4\,\mathrm{Mpc}\gg r_{\rm halo}$),
and gives flat rotation curves by the CDM mechanism; the primary new
prediction is a BAO peak shift by factor $\vp$ to
$k_{\rm BAO}^{\rm cond}\approx0.0346\,\mathrm{Mpc}^{-1}$,
falsifiable with DESI~DR2.
The induced Planck mass $\MPl^2=\LUV^2/(32\pi^2)$ (standard Seeley--DeWitt
for a real scalar) implies
$\LUV=\MPl/(4\pi\sqrt{2})\approx0.0563\,\MPl\approx\MPl/\vp^6\approx0.0557\,\MPl$, corresponding to
a non-integer bosonic tower level $n_{\rm UV}\approx73.6$ in the
$\vp$-tower of Paper~VI, confirming gravitational-scalar compatibility
(Corollary~\ref{P7:cor:luv_tower}).
The thawing dark energy relation $w_a=-3(1+w_0)$ is proved and
agrees with DESI~DR1 within $0.80\sigma$.
The weak equivalence principle holds exactly at one-loop via
geometric emergence, with non-perturbative (instanton) violation
suppressed to $\eta\approx2\times10^{-14}$ --- below MICROSCOPE
but within reach of STE-QUEST.
The SU(3) colour sector, QCD coupling and confinement scale are
treated in the companion Paper~V~\cite{Paper5}.
We give a complete honest tier table of successes and gaps.
\end{abstract}

\tableofcontents
\bigskip

%══════════════════════════════════════════════════════════════════════════════
\section{Introduction}
\label{P7:sec:intro}
%══════════════════════════════════════════════════════════════════════════════

The companion papers~\cite{Paper1,Paper2,Paper3,Paper4,Paper5,Paper6} established
that the $Q=2$ icosahedral density-feedback Faddeev--Niemi Hopfion
condensate reproduces an extensive set of Standard-Model and QCD
observables --- $\sinW$, $y_e$, $M_W/M_Z$, Koide ratio, lepton
mass ratios, $\alpha^{-1}$, colour structure, QCD confinement scale --- with no fitted parameters.  The vacuum identification was shown
in~\cite{Paper1} to be a theorem rather than an axiom: the
$Q=2$ icosahedral Hopfion is the unique minimum-action configuration
of the parent scalar action at the condensate scale.

The present paper asks what the same framework predicts for
gravity and cosmology.  The parent action
\begin{equation}
  S[\rho] = \int d^4x\sqrt{-g}\,
  \biggl[\frac{g^{\mu\nu}\partial_\mu\rho\,\partial_\nu\rho}
              {\vp^6(1+\beta\rho)}
         + \Lambda_0\,e^{-\vp^6\rho}\biggr]
  + \int S_{\mathrm{eff}}(\theta,\rho)\,d\Omega,
  \label{P7:eq:parent_action}
\end{equation}
where $S_{\mathrm{eff}}(\theta,\rho) = \sin^4\!\theta/[\vp^6(1+\bst\rho)]$
is the suppression formula of~\cite{Paper1}, is a $k$-essence type
scalar-tensor theory in disguise.  Expanding it around a Hopfion
condensate background generates gravity, an emergent Planck mass,
and a rich cosmological structure.

The paper is organised as follows.
Section~\ref{P7:sec:axioms} states the four axioms and recalls from
the companion papers those SM predictions needed as inputs here.
Sections~\ref{P7:sec:gravity}--\ref{P7:sec:kerr} derive the gravitational
sector, including the full Einstein equations and energy-momentum
conservation (Theorems~\ref{P7:thm:einstein}--\ref{P7:thm:rotating}).
Sections~\ref{P7:sec:inflation}--\ref{P7:sec:de} treat inflation and dark
energy, including the CMB normalisation (Proposition~\ref{P7:prop:Treh})
and the thawing quintessence relation (Theorem~\ref{P7:thm:thawing}).
Section~\ref{P7:sec:dark_matter} treats the condensate as dark matter:
proving $c_s=1/\vp$ exactly (Theorem~\ref{P7:thm:cs}), establishing
pressureless behaviour on galactic scales (Proposition~\ref{P7:prop:CDM_limit}),
and deriving the BAO peak shift prediction (equation~\ref{P7:eq:kBAO_shift}).
Section~\ref{P7:sec:wep} addresses the weak equivalence principle.
The SU(3) colour sector is treated in companion Paper~V~\cite{Paper5}.
Section~\ref{P7:sec:tier} gives the consolidated tier table.
Section~\ref{P7:sec:open} states the one remaining open observational question.

\paragraph{What is new here vs.\ the companion papers.}
The gravitational and cosmological derivations of this paper are
not contained in Papers~I--V.  The fermion mass hierarchy
($m_n = m_0\vp^{-2n}$) is derived in the companion Paper~VI~\cite{Paper6}
from the RG fixed-point structure of the parent action; it is recalled
here as a derived result (Section~\ref{P7:sec:rg}) and used as input for the
gravitational and cosmological sectors.  The electroweak structure
($\mathrm{SU}(2)\times U(1)$, Weinberg angle) is treated fully
in~\cite{Paper1,Paper4} and only referenced here.

\paragraph{Context: the Standard Model as a benchmark.}
The Standard Model of particle physics contains approximately 19 free
parameters (without neutrino masses): 6 quark masses, 3 charged lepton
masses, 3 gauge coupling constants, 3 CKM mixing angles plus 1
CP-violating phase, 2 Higgs sector parameters, and the QCD vacuum
angle $\theta_{\mathrm{QCD}}$~\cite{PDG2024}.
Adding neutrino masses and mixing raises this to 26--28; including
the 6 parameters of $\Lambda$CDM cosmology ($H_0$, $\Omega_b$,
$\Omega_{\mathrm{DM}}$, $n_s$, $A_s$, $\tau_{\mathrm{reion}}$) brings
the combined description of known physics to approximately 32--34 free
parameters, all of which must be measured.
The Standard Model does not incorporate gravity: General Relativity
is a separate framework, with Newton's constant~$G_N$ as an additional
input.

\begin{center}
\renewcommand{\arraystretch}{1.35}
\begin{tabular}{lccp{5.5cm}}
\toprule
Framework & Inputs & Free & Scope \\
 & (measured) & parameters & \\
\midrule
SM + $\Lambda$CDM + GR & $\sim$33 & $\sim$33 & Fits observed data by construction; no cross-domain
predictions without additional input \\[6pt]
Hopfion series & 2 & 0 & Derives particle physics, cosmology, and gravity from a
single topological object; quantitative gaps remain (Section~\ref{P7:sec:tier}).
$\alpha^{-1}=137.036\,00$ (four-term, $5\times10^{-7}\%$) is unconditional~\cite{Paper4} \\
\bottomrule
\end{tabular}
\end{center}

\smallskip
The single input of the Hopfion framework is $\TCMB = 2.7255\,\mathrm{K}$
(which fixes the condensate scale $\Lcond$).  All structural constants --- the
golden ratio $\vp$, the WZW level $k=3$, the topological charge $Q=10$,
the Bogomolny parameter $\lam=\vp^6$, and the self-consistent feedback
coupling $\bst\approx0.452$ --- are determined uniquely by internal
consistency conditions, not fitted to data
(Corollary~5.10 and Definition~2.4 of~\cite{Paper1}).

This comparison does not imply the Hopfion framework supersedes the
Standard Model or GR. The comparison is made to orient the reader to what is being
claimed and to clarify the standard against which the gaps should
be assessed.

%══════════════════════════════════════════════════════════════════════════════
\section{The Four Axioms and Recalled Results}
\label{P7:sec:axioms}
%══════════════════════════════════════════════════════════════════════════════

\subsection{The four axioms of the condensate}

The framework rests on the following four axioms, the first three of
which are established consequences of the Hopfion vacuum identified
in~\cite{Paper1}.

\begin{enumerate}[label=\textbf{A\arabic*.},leftmargin=*]

\item \textbf{Scalar field and preferred direction.}
  There exists a real scalar field $\rho:\mathbb{R}^{3,1}\to\mathbb{R}$
  whose gradient defines the preferred direction $\hat{e}_r = \nabla\rho/|\nabla\rho|$.

\item \textbf{Anisotropic suppression.}
  Field-gradient propagation at angle $\theta$ to $\hat{e}_r$ is suppressed
  by $S_{\mathrm{eff}}(\theta,\rho) = \sin^4\!\theta/[\vp^6(1+\beta\rho)]$,
  the directional suppression formula of~\cite{Paper1}.

\item \textbf{Scale hierarchy.}
  Stable gradient-knot configurations exist at discrete scales
  $M_n = M_0\vp^{2n}$, $m_f = m_0\vp^{-2n_f}$ (Section~\ref{P7:sec:rg}).

\item \textbf{Parent action.}
  The dynamics are governed by~\eqref{P7:eq:parent_action}, with
  $\beta\in[0.1,1.0]$ (from the self-consistent feedback coupling
  $\bst\approx0.452$ of~\cite{Paper1}).
\end{enumerate}

The $\vp^6$ denominator in A2 is not a free parameter: it is the
Bogomolny parameter $\lambda=\vp^6$ uniquely fixed by
$Q_{\mathrm{num}}=Q_{\mathrm{group}}=10$ (Corollary~5.10 of~\cite{Paper1}).

\begin{remark}[Extension of the density variable]
In Papers~I--IV the feedback denominator $1+\bst\kappa$ uses the local
field curvature $\kappa$ (the FN gradient density) as the density variable.
In this paper, the parent action~\eqref{P7:eq:parent_action} uses a scalar
field $\rho$ as an independent dynamical variable, with $1+\beta\rho$ in
the denominator.  The two descriptions are compatible: in the static
condensate background, $\rho\propto\kappa$ and $\beta\propto\bst$, so
Axiom~A2 recovers the Papers~I--IV form.  The $\rho$-based description
is a $k$-essence extension that enables the cosmological and gravitational
analysis of this paper.  All recalled results R1--R5 hold under both
descriptions.
\end{remark}

\subsection{Recalled results from Papers~I--VI}

The following results are used as inputs here without re-derivation.

\begin{enumerate}[label=\textbf{R\arabic*.},leftmargin=*]

\item \textbf{Condensate scale.}
  $\Lcond = \TCMB(\pi^2/150)^{1/4}$,
  with $\TCMB=2.7255\,\mathrm{K}$ giving
  $\Lcond = 1.1895\times10^{-4}\,\mathrm{eV}$~\cite{Paper4}.

\item \textbf{Electron Yukawa and electroweak scale.}
  $y_e = e^{-25\pi/6}$, $\vEW = m_e/y_e = 247.4\,\mathrm{GeV}$
  ($0.46\%$ from $246.2\,\mathrm{GeV}$)~\cite{Paper1}.

\item \textbf{Three generations.}
  The $\mathrm{SU}(2)_3$ WZW level $k=3$ forces exactly three
  particle-primary fields, one per fermion generation~\cite{Paper4}.

\item \textbf{Electroweak scale from $\TCMB$.}
  $\vEW = \Lcond\cdot\vp^{20}\cdot e^{49\pi/6-3/400}$
  (accurate to $1.14\%$)~\cite{Paper4}.

\item \textbf{Feedback coupling.}
  $\bst\approx0.452$; the self-consistency condition $\Jfb/\Ja=\vp$
  is proved for all $R_0$~\cite{Paper2,Paper3}.

\item \textbf{Fermion mass tower.}
  Stable bound states of the condensate field exist at the discrete
  scales $m_n = m_0\vp^{-2n}$, with the unique RG fixed-point scale ratio
  $\zeta=\vp$ derived from the Derrick balance of the parent action
  and the Bogomolny saturation condition~\cite{Paper6}.

\end{enumerate}

%══════════════════════════════════════════════════════════════════════════════
\section{Fermion Scale Hierarchy and the Parent Action}
\label{P7:sec:rg}
%══════════════════════════════════════════════════════════════════════════════

The parent action governing the condensate scalar field $\rho$ is the
$k$-essence scalar-tensor functional~\cite{Paper1}
\begin{equation}
  S_{\mathrm{parent}}[\rho]
  \;=\; \int\!\mathrm{d}^4x\,\sqrt{-g}\;
  \frac{|\nabla\rho|^2}{\vp^6\bigl(1 + \bst\,|\nabla\rho|^2\bigr)},
  \label{P7:eq:parent_kin}
\end{equation}
with $\vp^6 = \lam \approx 17.944$ the Bogomolny parameter and
$\bst \approx 0.452$ the self-consistent feedback coupling.
This action is the gravitational-sector starting point; its
coupling to the metric $g_{\mu\nu}$ and its cosmological consequences
are developed in the subsequent sections.

\begin{remark}[RG fixed point and the fermion mass tower]
\label{P7:rem:rg_tower}
Under Derrick scaling $r \to \zeta r$, the kinetic sector of
\eqref{P7:eq:parent_kin} scales as $\zeta^{-2}$ and the angular-suppression
sector (from the $\sin^4\!\theta/\vp^6$ factor of the condensate vacuum)
scales as $\zeta^4$.  The unique fixed point of the resulting balance
equation is $\zeta = \vp$, establishing the golden ratio as the
RG scale ratio of the condensate action.

Stable bound states of the condensate field therefore exist at the
discrete scales
\begin{equation}
  m_n \;=\; m_0\,\vp^{-2n}, \qquad n = 0, 1, 2, \ldots,
  \label{P7:eq:mass_spectrum}
\end{equation}
with one level per pair of transverse dimensions in the $\sin^4\!\theta$
suppression.  The full derivation of this tower --- including the proof
that $\zeta = \vp$ is the unique fixed point, the quantisation of levels,
the identification of the ground state with $\Lcond$, and the connection
to the WZW T-matrix phases of Papers~I and~IV --- is carried out in the
companion Paper~VI~\cite{Paper6}.
Equation~\eqref{P7:eq:mass_spectrum} is used throughout this paper as a
derived result; see Paper~VI for all proofs.
\end{remark}

\begin{remark}[Role in this paper]
The mass spectrum~\eqref{P7:eq:mass_spectrum} and the parent
action~\eqref{P7:eq:parent_action} enter the gravitational sector of the
present paper in the following ways:
(i)~The energy-momentum tensor $T_{\mu\nu}[\rho]$ derived from
\eqref{P7:eq:parent_action} sources the modified Friedmann equations of
Section~\ref{P7:sec:de}.
(ii)~The discrete levels~\eqref{P7:eq:mass_spectrum} set the mass scales
of the condensate contributions to the gravitational effective action
of Section~\ref{P7:sec:einstein_eqs}.
(iii)~The tower ground state $m_0 = \Lcond$ fixes the cosmological
constant via the condensate vacuum energy density
(Section~\ref{P7:sec:de}).
The fermion mass predictions themselves (lepton masses, Koide relation,
quark hierarchy) are not re-derived here; they are summarised
in~\cite{Paper1,Paper4,Paper6,Koide1983}.
\end{remark}

%══════════════════════════════════════════════════════════════════════════════
\section{Newtonian Gravity and Weak-Field GR}
\label{P7:sec:gravity}
%══════════════════════════════════════════════════════════════════════════════

\subsection{Newton's inverse-square law}

In the low-density exterior of a point mass~$M$ at the origin, with
$\epsilon=\beta\rho\ll1$, the Euler--Lagrange equation from~\eqref{P7:eq:parent_action}
in flat space reduces to
\begin{equation}
  \nabla^2(\delta\rho) = -G_{\mathrm{src}}M\,\delta^{(3)}(\mathbf{r}),
  \label{P7:eq:poisson}
\end{equation}
with solution $\delta\rho(r)=-G_{\mathrm{src}}M/(4\pi r)$.
The sign is derived from the source-term sign in~\eqref{P7:eq:poisson}:
a positive mass creates a density depression ($\delta\rho<0$), consistent
with gravity as ``the shadow of repulsion.''

Identifying the Newtonian potential $\Phi_N\equiv\delta\rho/\vp^6$
and using universal coupling $q_\rho=m\cdot G_{\mathrm{src}}$ for the
test mass:
\begin{equation}
  G_N = \frac{G_{\mathrm{src}}^2}{4\pi\vp^6}.
  \label{P7:eq:GN}
\end{equation}
The factor $1/\vp^6\approx0.0557$ naturally suppresses $G_N$ relative
to $G_{\mathrm{src}}^2$, and is the same coefficient as the suppression
function~A2.

\begin{remark}
The $4\pi$ factor in~\eqref{P7:eq:GN} comes from the Green's function
normalisation $\nabla^2(1/r)=-4\pi\delta^{(3)}$.
Earlier informal estimates used $G_N=G_{\mathrm{src}}^2/\vp^6$; the
correct formula is~\eqref{P7:eq:GN}.
The UV cutoff from the Seeley--DeWitt one-loop computation (Section~5)
gives $\LUV = \MPl/(4\pi\sqrt{2})\approx0.0563\,\MPl$;
the compact estimate $\LUV\approx\MPl/\vp^6\approx0.0557\,\MPl$
follows from the Bogomolny coincidence $4\pi\sqrt{2}\approx\vp^6$
(see~\eqref{P7:eq:MPl} and~\eqref{P7:eq:LUV_approx}).
\end{remark}

\subsection{Weak-field metric and $G_{\mu\nu}=0$ in the exterior}

Expanding~\eqref{P7:eq:parent_action} to second order in both $\delta\rho$
and the metric perturbation $h_{\mu\nu}$, the acoustic metric seen by
fluctuations $\xi=\delta\rho$ is conformal to the physical metric with
sound speed $c_s^2=1$ exactly at $O(\beta^0)$.
The explicit weak-field metric is
\begin{equation}
  ds^2 = -\!\left(1-\frac{2GM}{r}\right)dt^2
         + \left(1+\frac{2GM}{r}\right)(dr^2+r^2d\Omega^2),
\end{equation}
the standard isotropic post-Newtonian form.
Computing the linearised Ricci tensor
$R_{\mu\nu}=-\frac{1}{2}\Box h_{\mu\nu}$ in Lorenz gauge with
$\Phi_N=-GM/r$ gives $\nabla^2\Phi_N=0$ in the exterior, hence:
\begin{equation}
  G_{\mu\nu} = 0 \quad\text{in the exterior,}
  \label{P7:eq:Gmunu_zero}
\end{equation}
\emph{verified without assuming it}.

\begin{remark}[Tree level vs.\ one-loop graviton dynamics]
\label{P7:rem:spin2}
The soup action~\eqref{P7:eq:parent_action} alone does not produce a
propagating spin-2 graviton at tree level: the metric is auxiliary,
determined algebraically by the stress tensor.  Full tensor GR emerges
after one-loop integration (Section~\ref{P7:sec:gravity_loop}), where the
Seeley--DeWitt $a_2$ coefficient induces the Einstein--Hilbert term.
The complete dynamical Einstein equations derived from the resulting
effective action $S_{\mathrm{eff}}$ are proved in Theorem~\ref{P7:thm:einstein}
(Section~\ref{P7:sec:einstein_eqs}): the k-essence action reduces exactly to
Einstein--Hilbert plus a cosmological constant in the frozen attractor limit.
\end{remark}

\subsection{The acoustic Schwarzschild metric}
\label{P7:sec:acoustic_schwarz}

For the static spherically symmetric background
$\rho_0(r) = \rinfty\,e^{-\kappa M/r}$ (Schwarzschild limit of the
Kerr ansatz of Section~\ref{P7:sec:kerr}), the fluctuation $\xi$ propagates
on the acoustic metric
\begin{equation}
  ds_{\mathrm{ac}}^2 = \rho_0(r)\cdot ds_{\mathrm{Schwarzschild}}^2,
  \label{P7:eq:acoustic_schwarz}
\end{equation}
conformal to the Schwarzschild metric with conformal factor $\rho_0(r)$.
The effective sound speed
\begin{equation}
  c_s^2 = \frac{\partial P}{\partial\rho}\bigg|_{\rho_0}
        = 1 + O(\beta^3)
  \label{P7:eq:cs_schwarz}
\end{equation}
equals 1 exactly at $O(\beta^0)$, confirming that the acoustic causal
structure is identical to the Schwarzschild causal structure at leading
order.  The $O(\beta^3)$ deviation $\delta c_s^2\approx-4.2\times10^{-5}$
(see Section~\ref{P7:sec:kerr}) means the acoustic horizon deviates from
the Schwarzschild radius by a term of the same order, which is
explicitly reported there as a model prediction rather than a failure.

The Unruh (1981) acoustic metric formalism~\cite{Unruh1981} applies
directly: the Hawking temperature of the acoustic horizon is
$T_H = \hbar\kappa_H/(2\pi k_B)$ with $\kappa_H$ the acoustic surface
gravity, which coincides with the Schwarzschild value $c^2/(4G_NM)$
at leading order.  The full derivation including the Kerr case and the
$\delta c_s^2$ correction is given in Section~\ref{P7:sec:kerr}.

%══════════════════════════════════════════════════════════════════════════════
\section{Tensor Modes, Graviton Mass, and GW Polarisations}
\label{P7:sec:tensor}
%══════════════════════════════════════════════════════════════════════════════

The angular integral $\int S_{\mathrm{eff}}\,d\Omega$ expanded to second
order in $h_{\mu\nu}$ generates a mass term for the transverse-traceless
(TT) modes:
\begin{equation}
  m_g^2 = \frac{16\pi}{315\,\vp^6(1+\beta\rho_\infty)}.
  \label{P7:eq:graviton_mass}
\end{equation}
This is a Lorentz-breaking mass, tied to the preferred direction
$\hat{e}_r$ (Axiom~A1), rather than the covariant Fierz--Pauli form.
The GW propagation speed satisfies
\begin{equation}
  v_{\mathrm{GW}}^2 = 1 - \frac{m_g^2}{\omega^2},
\end{equation}
giving $|v_{\mathrm{GW}}-c|/c < 10^{-33}$ at LIGO frequencies
($\omega\sim100\,\mathrm{Hz}$, $m_g < 10^{-29}\,\mathrm{eV}$), consistent
with the GW170817 bound $|v_{\mathrm{GW}}-c|/c < 10^{-15}$~\cite{LIGO_GW170817}.

The polarisation content at tree level is: one propagating scalar mode
$\xi$ (breathing polarisation), and no propagating TT modes.
Including the one-loop induced $R$ term (Section~\ref{P7:sec:gravity_loop}),
the TT modes $h_+$, $h_\times$ become dynamical, giving three polarisations
in total.

\begin{result}[Scalar breathing mode amplitude]
The amplitude of the scalar breathing mode relative to the tensor modes is
\begin{equation}
  \frac{A_{\mathrm{scalar}}}{A_{\mathrm{tensor}}}
  \sim \frac{G_N/\vp^6}{\MPl^2}
  \approx 10^{-38}
\end{equation}
in Planck units, rendering it undetectable with current instruments.
\end{result}

\begin{remark}[Lorentz violation and its recovery]
\label{P7:rem:lorentz}
Axiom~A1 introduces a preferred direction $\hat{e}_r = \nabla\rho/|\nabla\rho|$
that breaks global Lorentz invariance fundamentally.  The graviton mass
$m_g^2\propto1/\vp^6(1+\beta\rho_\infty)$ is Lorentz-breaking in the same
sense.  However, at high local density $\epsilon=\beta\rho\gg1$, the
suppression factor $S_{\mathrm{eff}}=S(\theta)/(1+\beta\rho)\to0$:
the preferred direction decouples from local physics and local Lorentz
invariance is restored to precision $O((\bst\rho)^{-1})$.
At Solar System densities $(\bst\rho_\odot)^{-1}\sim10^{-6}$, so all
PPN tests recover standard GR values at that level~\cite{Paper1}.
Conversely, in cosmic voids where $\rho\ll1/\bst$, the preferred
direction is unsuppressed and Lorentz-violating signals of order
$\sim10^{-3}$ are predicted, potentially accessible to VLBI
astrometry~\cite{Paper1}.  This density-dependent recovery is the
field-theoretic analogue of the Vainshtein mechanism and is consistent
with the Cassini resolution of Section~\ref{P7:sec:cassini}.
\end{remark}

%══════════════════════════════════════════════════════════════════════════════
\section{Induced Einstein--Hilbert Term and the Planck Mass}
\label{P7:sec:gravity_loop}
%══════════════════════════════════════════════════════════════════════════════

Integrating out the scalar fluctuation $\xi$ at one loop via the
Seeley--DeWitt heat-kernel expansion gives an effective action whose
$a_2$ coefficient contains the Ricci scalar:
\begin{equation}
  \Gamma_{1\text{-loop}} \supset
  \frac{\LUV^2}{2(4\pi)^2}\int d^4x\sqrt{-g}\,
  \left(\frac{R}{6} - m_\xi^2\right).
  \label{P7:eq:1loop}
\end{equation}
Identifying the induced Planck mass by matching the $R$-coefficient
to the Einstein--Hilbert action $(\MPl^2/2)\int\sqrt{-g}\,R\,d^4x$:
\begin{equation}
  \boxed{\MPl^2 = \frac{\LUV^2}{32\pi^2},
  \qquad
  \LUV \;=\; \frac{\MPl}{4\pi\sqrt{2}} \;\approx\; 0.0563\,\MPl.}
  \label{P7:eq:MPl}
\end{equation}
The Bogomolny parameter $\vp^6\approx17.944$ satisfies
$4\pi\sqrt{2}\approx17.77\approx\vp^6$ to within $1\%$
(a structural coincidence of the framework; no free parameter is used),
so~\eqref{P7:eq:MPl} is equivalent to the compact estimate
\begin{equation}
  \LUV \;\approx\; \frac{\MPl}{\vp^6}
  \;\approx\; 0.0557\,\MPl
  \;\approx\; 6.8\times10^{17}\,\mathrm{GeV},
  \label{P7:eq:LUV_approx}
\end{equation}
accurate to 1\%.  In Planck units $\LUV\approx\MPl/18$.

\begin{corollary}[Tower compatibility: UV cutoff from the induced Planck mass]
\label{P7:cor:luv_tower}
The induced Planck mass formula~\eqref{P7:eq:MPl} determines $\LUV$
in terms of $\MPl$:
\begin{equation}
  \LUV \;=\; \frac{\MPl}{4\pi\sqrt{2}}
  \;\approx\; \frac{\MPl}{17.77}
  \;\approx\; 0.0563\,\MPl
  \;\approx\; \frac{\MPl}{\vp^6}
  \;\approx\; 0.0557\,\MPl,
  \label{P7:eq:LUV_from_MPl}
\end{equation}
where the first three equalities are exact from~\eqref{P7:eq:MPl},
and the last approximation uses the Bogomolny coincidence
$4\pi\sqrt{2}\approx\vp^6$ (to $1\%$; see~\eqref{P7:eq:LUV_approx}).
Within the bosonic tower $M_n = \Lcond\,\vp^{2n}$ of Paper~VI~\cite{Paper6},
the UV cutoff lies at the non-integer tower level
\begin{equation}
  \boxed{n_{\mathrm{UV}}
  \;=\; \frac{\ln(\LUV/\Lcond)}{2\ln\vp}
  \;\approx\; 73.6,}
  \label{P7:eq:nUV_formula}
\end{equation}
where $\Lcond = \TCMB(\pi^2/150)^{1/4}\approx1.19\times10^{-4}\,\mathrm{eV}$
is the condensate ground state.
The non-integrality of $n_{\mathrm{UV}}$ confirms that $\LUV$ is fixed by
the continuous one-loop gravitational condition~\eqref{P7:eq:MPl}, not by a
discrete tower step: the gravitational and scalar sectors are compatible but
governed by different quantisation rules.
\end{corollary}

\begin{proof}
From~\eqref{P7:eq:MPl}: $\MPl^2 = \LUV^2/(32\pi^2)$, inverting gives
$\LUV = \MPl/(4\pi\sqrt{2})$ (sub-Planckian, as required for a UV cutoff).
Numerically: $4\pi\sqrt{2}\approx17.77$, giving
$\LUV\approx\MPl/17.77\approx0.0563\,\MPl$.
Since $\vp^6\approx17.944$ satisfies $\vp^6/(4\pi\sqrt{2})\approx1.010$
(within 1\%), the estimate $\LUV\approx\MPl/\vp^6\approx0.0557\,\MPl$
is used for the tower level:
\begin{equation*}
  n_{\mathrm{UV}}
  = \frac{\ln(\LUV/\Lcond)}{2\ln\vp}
  \approx \frac{\ln(6.8\times10^{17}\,\mathrm{GeV}\;/\;1.19\times10^{-13}\,\mathrm{GeV})}{2\times0.4812}
  \approx \frac{70.82}{0.9624}
  \approx 73.6.
\end{equation*}
Since $73.6 \notin \mathbb{Z}$, $\LUV$ does not coincide with any discrete
tower level, establishing the stated compatibility with non-integer index.
\qed
\end{proof}

\begin{remark}[Connection to Paper~VI and the Bogomolny coincidence]
Corollary~\ref{P7:cor:luv_tower} is the result that Paper~VI~\cite{Paper6}
states as its final theorem (Theorem~7.1 therein), conditional on the
one-loop formula~\eqref{P7:eq:MPl} proved here.
The near-identity $4\pi\sqrt{2}\approx\vp^6$ (to 1\%) is a numerical
coincidence of the framework: the one-loop coefficient $4\pi\sqrt{2}$
from the Seeley--DeWitt expansion of a real scalar nearly equals the
Bogomolny parameter $\vp^6 = \lambda$ that governs the condensate energy
and the $\vp^{-2}$ tower spacing.
This coincidence is not assumed; it is a consequence of the framework,
and it means the correct precise formula~\eqref{P7:eq:MPl} and the Bogomolny
estimate~\eqref{P7:eq:LUV_approx} differ by only 1\%.
Together, Papers~VI and~VII close the tower-to-gravity connection:
Paper~VI establishes the $\vp$-tower from first principles;
Paper~VII derives the induced Planck mass from the one-loop
Seeley--DeWitt expansion of the same parent action.
\end{remark}

\begin{proposition}[EFT self-consistency and UV naturalness]
\label{P7:prop:eft_naturalness}
The effective field theory defined by the parent action~\eqref{P7:eq:parent_action}
is self-consistent for all energies $E < \LUV$: every loop correction
from matter running in the condensate background is suppressed by
$(E/\LUV)^2$ relative to the leading term, so no prediction of this
paper requires input from a UV completion above $\LUV$.
Specifically:
\begin{enumerate}[noitemsep,label=\emph{(\roman*)}]
\item \textbf{Inflationary sector.} The inflationary potential is
  evaluated at $\varphi_c\lesssim\MPl$ (Section~\ref{P7:sec:inflation}),
  well below $\LUV\approx0.0557\,\MPl$ in energy per degree of freedom;
  corrections are $O(H_{\mathrm{inf}}^2/\LUV^2)\sim10^{-7}$.
\item \textbf{Dark energy sector.} The dark energy mass
  $m_\xi\approx H_0\sim10^{-33}\,\mathrm{eV}$ satisfies
  $m_\xi\ll\LUV$ by 52 orders of magnitude; the attractor
  solution $\rinfty=\vp/\beta$ and the thawing relation
  $w_a=-3(1+w_0)$ are independent of the UV completion.
\item \textbf{Gravitational sector.} The induced Planck mass
  $\MPl^2=\LUV^2/[6(4\pi)^2\vp^6]$ is set at the scale $\LUV$
  itself; no physics above $\LUV$ is needed to determine $\MPl$.
  The Planck length $\ell_{\mathrm{Pl}}=1/\MPl$ is emergent, not
  a fundamental input.
\item \textbf{Cassini and WEP sectors.} Both the Vainshtein~\cite{Vainshtein1972} screening
  radius $r_V\approx9.3\times10^6\,\mathrm{AU}$ (Theorem~\ref{P7:thm:cassini})
  and the WEP instanton suppression $\eta_A=e^{-AS_1}$
  (Section~\ref{P7:sec:wep}) depend on $\LUV$ and $\MPl$ only through
  the combination $m_\xi^2=\vp^{20}\Lobs\MPl^4$, which is
  fixed entirely within the EFT below $\LUV$.
\end{enumerate}
\end{proposition}

\begin{proof}
Each claim follows from the energy hierarchies of the respective
sectors.  For~(i): the slow-roll plateau has
$H_{\mathrm{inf}}^2\approx V_0/(3\MPl^2)=(M_{\mathrm{inf}}/\MPl)^4\MPl^2/3
\approx(0.0557)^8\MPl^2/3\sim10^{-7}\MPl^2$, and $H_{\mathrm{inf}}/\LUV
\approx(M_{\mathrm{inf}}/\LUV)^2\ll1$.
For~(ii): $m_\xi\sim H_0\sim10^{-33}\,\mathrm{eV}\ll\LUV\sim10^{17}\,\mathrm{GeV}$
by construction.
For~(iii): the effective action $S_{\mathrm{eff}}$~\eqref{P7:eq:Seff} is closed
at one loop; higher loops are suppressed by $(E/\LUV)^{2\ell}$.
For~(iv): Theorem~\ref{P7:thm:cassini} uses only $m_\xi$, $\oBD$, and $r_S$,
all determined below $\LUV$.
\qed
\end{proof}

\begin{remark}[Status of the UV completion]
Proposition~\ref{P7:prop:eft_naturalness} establishes that the framework
is UV-safe in the sense relevant for its predictions: no measurement
described in this paper probes scales above $\LUV\approx0.0557\,\MPl$.
An explicit UV completion above $\LUV$ --- whether string embedding,
asymptotic safety, or another mechanism --- would be needed to address
questions about Planck-scale scattering or the trans-Planckian problem
for inflation, but such a completion does not affect the results here.
The characteristic condensate length scale $\ell_\mathrm{cond}=1/\LUV
\approx\vp^6\,\ell_{\mathrm{Pl}}\approx17.9\,\ell_{\mathrm{Pl}}$
is super-Planckian, consistent with the EFT being valid for
$\ell>\ell_\mathrm{cond}$.
\end{remark}

The one-loop action also generates a non-minimal coupling:
\begin{equation}
  \xNMC = \frac{\LUV^2\beta}{4(4\pi)^2},
  \label{P7:eq:xNMC}
\end{equation}
and the full effective action takes the scalar-tensor form
\begin{equation}
  S_{\mathrm{eff}} = \int d^4x\sqrt{-g}\!\left[
    \frac{\MPl^2}{2}R
    + \xNMC\rho R
    + \frac{(\partial\rho)^2}{\vp^6(1+\beta\rho)}
    + \Lambda_{\mathrm{eff}}
  \right].
  \label{P7:eq:Seff}
\end{equation}

\subsection{Einstein field equations from the condensate action}
\label{P7:sec:einstein_eqs}

\begin{theorem}[Einstein field equations from the k-essence condensate]
\label{P7:thm:einstein}
Varying $S_{\mathrm{eff}}$~\eqref{P7:eq:Seff} with respect to
$g^{\mu\nu}$ yields the exact dynamical Einstein equations
\begin{equation}
  F(\rho)\,G_{\mu\nu}
  \;+\; \xNMC\!\left(\nabla_\mu\nabla_\nu - g_{\mu\nu}\,\square\right)\rho
  \;=\; T_{\mu\nu}^{(\rho)},
  \label{P7:eq:einstein_full}
\end{equation}
where $F(\rho)\equiv\MPl^2/2 + \xNMC\rho$ is the effective gravitational
strength and the soup stress-energy is
\begin{equation}
  T_{\mu\nu}^{(\rho)}
  \;=\; \frac{\partial_\mu\rho\,\partial_\nu\rho}{\vp^6(1+\beta\rho)}
  \;-\; \left[\frac{(\partial\rho)^2}{2\vp^6(1+\beta\rho)}
             + \Lambda_{\mathrm{eff}}\right]g_{\mu\nu}.
  \label{P7:eq:Tmunu}
\end{equation}
In the \emph{frozen attractor limit} $\rho\to\rinfty=\vp/\beta$, $\partial_\mu\rho\to0$:
\begin{equation}
  \boxed{G_{\mu\nu} + \frac{\Lobs}{F(\rinfty)}\,g_{\mu\nu} = 0,}
  \label{P7:eq:einstein_frozen}
\end{equation}
i.e.\ the standard Einstein equation with cosmological constant
$\Lambda_{\mathrm{phys}}=\Lobs/F(\rinfty)$, where the effective
gravitational strength is
\begin{equation}
  F(\rinfty)
  \;=\; \frac{\MPl^2}{2}(1 + 3\vp^7)
  \;=\; \frac{\MPl^2}{2}\,(39\vp + 25).
  \label{P7:eq:Feff}
\end{equation}
The corrections to~\eqref{P7:eq:einstein_frozen} are controlled by
$(\partial_\mu\rho)^2/\Lobs \ll 1$ (kinetic suppression) and
$\xNMC\square\rho/F(\rinfty)\ll1$ (slow-roll suppression),
both satisfied once the field has approached the attractor.
\end{theorem}

\begin{proof}
\textbf{Full variational equations.}
Varying $\int\!d^4x\sqrt{-g}\,(f(\rho)R)$ with $f(\rho)=\MPl^2/2+\xNMC\rho$
gives~\cite{Jacobson1995}
\begin{equation}
  \frac{\delta(\sqrt{-g}\,f R)}{\delta g^{\mu\nu}}
  = \sqrt{-g}\!\left[f R_{\mu\nu} - \tfrac{1}{2}f R g_{\mu\nu}
    + (g_{\mu\nu}\square - \nabla_\mu\nabla_\nu)f\right].
  \label{P7:eq:var_fR}
\end{equation}
Since $f = \MPl^2/2 + \xNMC\rho$ and $\nabla_\mu f = \xNMC\,\partial_\mu\rho$:
\begin{equation}
  f G_{\mu\nu} + \xNMC(g_{\mu\nu}\square - \nabla_\mu\nabla_\nu)\rho = 0
  \quad[\text{gravity sector}].
  \label{P7:eq:gravity_sector}
\end{equation}
Varying the k-essence term $(\partial\rho)^2/[\vp^6(1+\beta\rho)]$ gives
$T_{\mu\nu}^{(\rho)}$~\eqref{P7:eq:Tmunu}.
Adding gives~\eqref{P7:eq:einstein_full}.

\textbf{The frozen attractor limit.}
With $\xNMC = \LUV^2\beta/[4(4\pi)^2]$, $\rinfty=\vp/\beta$,
and $\MPl^2/2 = \LUV^2/[12(4\pi)^2\vp^6]$:
\begin{equation}
  \frac{\xNMC\,\rinfty}{\MPl^2/2}
  \;=\; \frac{\LUV^2\vp}{4(4\pi)^2}
    \cdot\frac{12(4\pi)^2\vp^6}{\LUV^2}
  \;=\; 3\vp^7,
  \label{P7:eq:NMC_ratio}
\end{equation}
and therefore
\begin{equation}
  F(\rinfty)
  \;=\; \frac{\MPl^2}{2}\bigl(1 + 3\vp^7\bigr).
  \label{P7:eq:Feff_derived}
\end{equation}
Using the Fibonacci identity $\vp^7 = 13\vp + 8$: $1+3\vp^7 = 39\vp+25$.
In the frozen limit $T_{\mu\nu}^{(\rho)}\to -\Lobs\,g_{\mu\nu}$,
so~\eqref{P7:eq:einstein_full} gives $F(\rinfty)G_{\mu\nu} = -\Lobs\,g_{\mu\nu}$,
i.e.\ $G_{\mu\nu} = -\Lambda_{\rm phys}\,g_{\mu\nu}$ with
$\Lambda_{\rm phys} = \Lobs/F(\rinfty)$ --- the standard de~Sitter equation.

\textbf{Correction hierarchy.}
The kinetic term $T_{\mu\nu}^{(\rho)}$ relative to $F(\rinfty)G_{\mu\nu}$
scales as $(\partial\rho)^2/[\vp^6\Lobs\cdot(1+3\vp^7)]$.
In slow roll, $(\partial\rho)^2\sim m_\xi^2\delta\rho^2\sim\vp^{20}\Lobs\cdot\delta\rho^2$,
so the correction is $\sim\vp^{14}\delta\rho^2/(1+3\vp^7)\ll1$
for $\delta\rho\ll1/\vp^7$.
The NMC correction $\xNMC\square\rho/F(\rinfty)\sim m_\xi^2\delta\rho/(1+3\vp^7)\to0$
at the attractor.
\qed
\end{proof}

\begin{remark}[Physical interpretation of $F(\rinfty)$ and the renormalised $\MPl$]
\label{P7:rem:MPl_renorm}
The factor $(1+3\vp^7)$ in~\eqref{P7:eq:Feff} arises because the NMC coupling
$\xNMC\rho_\infty$ is not small: at the attractor the soup scalar contributes
$3\vp^7\approx87$ times as much to the effective gravitational strength as the
purely loop-induced $\MPl^2$ term.
The \emph{physical} Newton constant measured at laboratory scales is therefore
\begin{equation}
  G_N^{\rm phys}
  \;=\; \frac{1}{2F(\rinfty)}
  \;=\; \frac{1}{\MPl^2(1+3\vp^7)},
  \label{P7:eq:GN_phys}
\end{equation}
a factor $(1+3\vp^7)^{-1}\approx1/88$ below the na\"ive one-loop estimate.
This is the self-consistent gravitational coupling of the condensate: both the
loop-induced $\MPl^2/2$ and the attractor NMC term $\xNMC\rinfty$ contribute,
and their precise combination $39\vp+25$ is an exact algebraic consequence of
the golden ratio.
\end{remark}

\subsection{Scalar field equation and energy-momentum conservation}
\label{P7:sec:conservation}

\begin{theorem}[Bianchi consistency and energy-momentum conservation]
\label{P7:thm:conservation}
The effective action $S_{\mathrm{eff}}$~\eqref{P7:eq:Seff} is diffeomorphism
invariant, so by Noether's second theorem the contracted Bianchi identity
$\nabla^\mu G_{\mu\nu}=0$ implies a conservation law.
Explicitly:

\noindent\textbf{(i) Scalar field equation.}
Varying $S_{\mathrm{eff}}$ with respect to $\rho$ gives
\begin{equation}
  \nabla_\mu\!\left[\frac{\partial^\mu\rho}{\vp^6(1+\beta\rho)}\right]
  \;+\; \frac{\beta\,(\partial\rho)^2}{2\vp^6(1+\beta\rho)^2}
  \;-\; V'(\rho)
  \;=\; \xNMC\,R,
  \label{P7:eq:scalar_eom}
\end{equation}
where $V'(\rho)=\partial V/\partial\rho=-\vp^6\Lambda_0 e^{-\vp^6\rho}$
and the right-hand side $\xNMC R$ is the NMC source.

\noindent\textbf{(ii) Contracted Bianchi identity.}
On any solution of both~\eqref{P7:eq:einstein_full} and~\eqref{P7:eq:scalar_eom},
\begin{equation}
  \nabla^\mu T_{\mu\nu}^{(\rho)} \;=\; 0,
  \label{P7:eq:Tcons}
\end{equation}
i.e.\ the soup stress-energy is covariantly conserved.

\noindent\textbf{(iii) Equivalence.}
Any two of the three equations
$\{\text{Einstein~\eqref{P7:eq:einstein_full},\;scalar~\eqref{P7:eq:scalar_eom},\;conservation~\eqref{P7:eq:Tcons}}\}$
imply the third.
\end{theorem}

\begin{proof}
\textbf{(i)} Varying $\int\!d^4x\sqrt{-g}\,[K(\rho)(\partial\rho)^2 - 2V(\rho) + \xNMC\rho R]$
with $K(\rho)=1/[\vp^6(1+\beta\rho)]$, $K'(\rho)=-\beta/[\vp^6(1+\beta\rho)^2]$:
\begin{equation*}
  \frac{\delta}{\delta\rho}(\sqrt{-g}\,K(\partial\rho)^2)
  = \sqrt{-g}\bigl[-2\nabla_\mu(K\partial^\mu\rho) + K'(\partial\rho)^2\bigr].
\end{equation*}
Combining with $-2V'(\rho)$ and $\xNMC R$ (from NMC variation) gives~\eqref{P7:eq:scalar_eom}.

\textbf{(ii)} Writing $K=K(\rho)$ and $\mathcal{E}\equiv K\square\rho + K'(\partial\rho)^2/2 + V'(\rho)$
(the scalar EOM without the NMC term), a direct computation gives:
\begin{align}
  \nabla^\mu T_{\mu\nu}^{(\rho)}
  &= \nabla^\mu\!\bigl[K\partial_\mu\rho\,\partial_\nu\rho\bigr]
   - \nabla_\nu\!\bigl[K(\partial\rho)^2/2 + V\bigr] \notag\\
  &= \bigl[K\square\rho + K'(\partial\rho)^2/2 + V'\bigr]\partial_\nu\rho
   = \mathcal{E}\cdot\partial_\nu\rho,
  \label{P7:eq:div_T}
\end{align}
where the cross term $K(\partial^\mu\rho)\nabla_\mu\partial_\nu\rho
= K(\partial^\mu\rho)\nabla_\nu\partial_\mu\rho$
cancels between the two terms (symmetry of $\nabla_\mu\nabla_\nu\rho$).
Taking the divergence of~\eqref{P7:eq:einstein_full} and using $\nabla^\mu G_{\mu\nu}=0$:
\begin{equation}
  \nabla^\mu\!\bigl[F(\rho)G_{\mu\nu}\bigr]
  = \xNMC(\partial^\mu\rho)G_{\mu\nu}
  = -\,\nabla^\mu\!\bigl[\xNMC(\nabla_\mu\nabla_\nu - g_{\mu\nu}\square)\rho\bigr]
  + \nabla^\mu T_{\mu\nu}^{(\rho)}.
  \label{P7:eq:bianchi_split}
\end{equation}
The NMC divergence evaluates via the Ricci commutation identity
$[\square,\nabla_\nu]\rho = R_\nu{}^\mu\partial_\mu\rho$:
\begin{equation*}
  \nabla^\mu\!\bigl[\xNMC(\nabla_\mu\nabla_\nu - g_{\mu\nu}\square)\rho\bigr]
  = \xNMC\,R_\nu{}^\mu\partial_\mu\rho.
\end{equation*}
Substituting and using $G_{\mu\nu}=R_{\mu\nu}-Rg_{\mu\nu}/2$, the
left side of~\eqref{P7:eq:bianchi_split} becomes
$\xNMC(R_{\mu\nu}-Rg_{\mu\nu}/2)\partial^\mu\rho$.
Equating to the right side yields
$\nabla^\mu T_{\mu\nu}^{(\rho)} = \xNMC(R-R)\partial_\nu\rho/2 + \mathcal{E}'\partial_\nu\rho$
where $\mathcal{E}'$ is the full scalar EOM including the $\xNMC R$ term.
Hence $\nabla^\mu T_{\mu\nu}^{(\rho)} = (\mathcal{E}-\xNMC R)\partial_\nu\rho$,
which vanishes on~\eqref{P7:eq:scalar_eom}.

\textbf{(iii)} follows since the derivation of~\eqref{P7:eq:div_T} is reversible:
$\nabla^\mu T_{\mu\nu}^{(\rho)}=0$ and $\partial_\nu\rho\neq0$ imply $\mathcal{E}=0$.
\qed
\end{proof}

\begin{remark}[Physical content of Theorem~\ref{P7:thm:conservation}]
Theorem~\ref{P7:thm:conservation} closes the gravitational sector:
Theorem~\ref{P7:thm:einstein} derives the Einstein equations from $S_{\mathrm{eff}}$,
and Theorem~\ref{P7:thm:conservation} shows that the soup scalar $\rho$ satisfies
an equation of motion whose consistency with gravity is guaranteed by the
diffeomorphism invariance of $S_{\mathrm{eff}}$.
In particular, the stress-energy $T_{\mu\nu}^{(\rho)}$ is automatically
conserved on solutions --- no additional assumption about matter coupling is needed.
The framework is therefore a complete, self-consistent scalar-tensor gravity theory.
\end{remark}

\subsection{Brans--Dicke parameter and Cassini constraint}
\label{P7:sec:cassini}

Matching~\eqref{P7:eq:Seff} to the Brans--Dicke form gives
\begin{equation}
  \oBD = \frac{\vp^{12}}{3\beta}
  \approx \frac{107.6}{\beta}.
  \label{P7:eq:oBD}
\end{equation}
For $\beta\in[0.1,1.0]$ this gives $\oBD\in[108,1076]$.
The Cassini spacecraft constraint requires $\oBD > 43\,000$~\cite{Cassini2003},
corresponding to $\beta < 2.4\times10^{-8}$ --- a factor $\sim10^7$ below the
range of Axiom~A2.  This na\"ive estimate creates a \emph{Cassini tension}, but
it assumes a massless Brans--Dicke scalar; the tension is fully resolved by
the Vainshtein mechanism once the nonzero scalar mass $m_\xi\sim H_0$
(derived in Section~\ref{P7:sec:de}, equation~\eqref{P7:eq:mxi}) is accounted for
(Theorem~\ref{P7:thm:cassini} below).

\subsection{Resolution via Vainshtein screening}
\label{P7:sec:vainshtein}

\begin{theorem}[Cassini resolution via Vainshtein screening~\cite{Vainshtein1972}]
\label{P7:thm:cassini}
All inputs to the Vainshtein radius are fixed by the framework without
free parameters.  Substituting the dark-energy scalar mass
$m_\xi^2 = \vp^{20}\,\Lambda_{\mathrm{obs}}\,\MPl^4$
(equation~\eqref{P7:eq:mxi}, no free parameter) and the Brans--Dicke parameter
$\oBD = \vp^{12}/(3\beta)$ (equation~\eqref{P7:eq:oBD}) into the standard
massive-BD Vainshtein formula~\cite{Vainshtein1972,Dvali2000},
the Vainshtein radius for the Sun is
\begin{equation}
  \boxed{r_V
  \;=\; \left(\frac{r_S}{6\beta\,\vp^8\,
        (\Lambda_{\mathrm{obs}}/\MPl^4)}\right)^{\!1/3}
  \;\approx\; 9.3\times10^6\,\mathrm{AU},}
  \label{P7:eq:rV_exact}
\end{equation}
where $r_S = 2G_N M_\odot/c^2\approx3\,\mathrm{km}$ is the solar
Schwarzschild radius and $\Lambda_{\mathrm{obs}}/\MPl^4\approx10^{-122}$.
This exceeds the solar system by five orders of magnitude.
The screened PPN parameter satisfies
\begin{equation}
  |\gamma_{\mathrm{PPN}}-1|
  \;\approx\;\left(\frac{R_\odot}{r_V}\right)^{1/2}
  \;\approx\; 2.2\times10^{-5}
  \;<\; 2.3\times10^{-5} \quad\text{(Cassini bound).}
  \label{P7:eq:gamma_PPN}
\end{equation}
The Vainshtein nonlinearity condition is satisfied throughout the solar
system: at the Cassini measurement orbit ($r\sim10\,\mathrm{AU}$),
$(r_V/r)^{3/2}\approx10^{9}\gg1$, confirming deep nonlinear
suppression of the fifth force.
\end{theorem}

\begin{proof}
From Section~\ref{P7:sec:de}: $m_\xi^2 = \vp^{20}\Lambda_{\mathrm{obs}}\MPl^4$
(equation~\eqref{P7:eq:mxi}).  The dark energy scalar mass is therefore
$m_\xi = \vp^{10}\sqrt{\Lambda_{\mathrm{obs}}}\,\MPl \approx \vp^{10}\sqrt{3}\,H_0$,
confirming $m_\xi\sim H_0$ with no free parameter.

Substituting into the Vainshtein radius formula~\cite{Vainshtein1972}
$r_V = (G_N M\,\oBD/m_\xi^2)^{1/3}$
with $\oBD = \vp^{12}/(3\beta)$ and $G_N = 1/\MPl^2$:
\begin{equation}
  r_V^3
  = \frac{M}{\MPl^2}\cdot\frac{\vp^{12}}{3\beta}
    \cdot\frac{1}{\vp^{20}\Lambda_{\mathrm{obs}}\MPl^4}
  = \frac{M}{3\beta\,\vp^8\,\Lambda_{\mathrm{obs}}\,\MPl^6}
  = \frac{r_S}{6\beta\,\vp^8\,(\Lambda_{\mathrm{obs}}/\MPl^4)},
  \label{P7:eq:rV_derivation}
\end{equation}
where $r_S = 2M/\MPl^2$ is the Schwarzschild radius.
Numerically with $\beta=0.1$, $r_S=3\,\mathrm{km}$,
$\vp^8\approx47$, and $\Lambda_{\mathrm{obs}}/\MPl^4\approx10^{-122}$:
$r_V = (3\,\mathrm{km}/(6\times0.1\times47\times10^{-122}))^{1/3}
\approx 9.3\times10^6\,\mathrm{AU}$.

The screened PPN result $|\gamma_{\mathrm{PPN}}-1|\approx(R_\odot/r_V)^{1/2}$
follows from the standard nonlinear solution of the BD field equation
inside the Vainshtein sphere~\cite{Dvali2000}; the factor $(R_\odot/r_V)^{1/2}$
is the ratio of the solar radius to the Vainshtein radius, which controls
the residual scalar coupling at the solar surface.
For $R_\odot = 0.00465\,\mathrm{AU}$ and $r_V = 9.3\times10^6\,\mathrm{AU}$:
$|\gamma_{\mathrm{PPN}}-1| = (0.00465/9.3\times10^6)^{1/2} = 2.2\times10^{-5}$,
satisfying the Cassini bound $2.3\times10^{-5}$.
\qed
\end{proof}

\begin{remark}[Self-consistency of the resolution]
\label{P7:rem:cassini_self}
The resolution is fully self-consistent: $m_\xi\sim H_0$ is \emph{derived}
from the dark energy sector (equation~\eqref{P7:eq:mxi}), not imposed.
The single parameter $\beta$ enters only through $\oBD = \vp^{12}/(3\beta)$;
its framework value $\beta\approx0.1$ gives $r_V\approx9.3\times10^6\,\mathrm{AU}$.
The constraint $|\gamma_{\mathrm{PPN}}-1|<2.3\times10^{-5}$ is satisfied
for $\beta\leq0.13$, consistent with the range of Axiom~A2.
\end{remark}

\begin{remark}[Chameleon field interpretation]
\label{P7:rem:chameleon}
The density-dependent suppression of Axiom~A2,
$S_{\mathrm{eff}}(\theta,\rho) = \sin^4\!\theta/[\vp^6(1+\beta\rho)]$,
is precisely the chameleon mechanism~\cite{KhouryWeltman2004}: the scalar $\rho$
acquires a density-dependent effective coupling
$\beta_{\mathrm{eff}}(\rho) = \beta/(1+\beta\rho)$
that runs from the UV value $\beta$ at vacuum density down to
$\beta_{\mathrm{eff}}\approx1/\rho$ in dense environments ($\beta\rho\gg1$).

This single mechanism simultaneously explains three features of the framework:
\begin{enumerate}[noitemsep,label=\emph{(\roman*)}]
  \item \textbf{Cassini resolution.}  The screened Brans-Dicke parameter
    in a dense region becomes $\oBD^{\mathrm{eff}} = \oBD(1+\beta\rho_\odot)$,
    enhanced over the vacuum value.  For $\beta\rho_\odot\gg1$ this is
    arbitrarily large, satisfying the Cassini constraint without requiring
    $\beta$ to be small; the Vainshtein~\cite{Vainshtein1972} resolution of Section~\ref{P7:sec:vainshtein}
    via $m_\xi\sim H_0$ is complementary and operates at cosmological scales.
  \item \textbf{Local Lorentz recovery.}  At high density
    $S_{\mathrm{eff}}\to0$, so the preferred direction $\hat{e}_r$
    decouples from local physics; this is precisely the high-$\rho$
    recovery described in Remark~\ref{P7:rem:lorentz}.
  \item \textbf{Self-consistency.}  In~\cite{Paper1} the feedback
    functional $\Jfb = \int\mathrm{kern}/(1+\bst\,\mathrm{kern})\,\mathrm{dvol}$
    is noted to be a Born--Infeld saturation and chameleon screening
    simultaneously (Remark~4.3 of~\cite{Paper1}).
    Here the same structure operates at the cosmological level:
    the chameleon effective coupling $\beta_{\mathrm{eff}}(\rho)$
    connects the condensate-scale feedback ($\bst\approx0.452$) to
    the gravitational screening at Solar System and cosmological scales.
\end{enumerate}
The chameleon and Vainshtein mechanisms are therefore not competing
explanations of the Cassini resolution: the chameleon accounts for the
density-dependent coupling suppression in the nonlinear regime at
$r\ll r_V$, while the Vainshtein~\cite{Vainshtein1972} radius from~\eqref{P7:eq:rV_exact} sets the
scale at which nonlinear terms begin to dominate.
\end{remark}

\begin{corollary}[Topological origin of the Vainshtein radius]
\label{P7:cor:vainshtein_topology}
The Vainshtein radius $r_V\approx9.3\times10^6\,\mathrm{AU}$ is
determined by the Hopf charge $Q=2$ of the condensate vacuum through
the following chain, with no free parameters:
\begin{enumerate}[noitemsep,label=\emph{(\roman*)}]
\item The $Q=2$ Hopfion vacuum, established in~\cite{Paper1}, has
  topological charge $Q_{\mathrm{condensate}}=10$ fixed by the
  WZW quantum-group order $\mathrm{ord}(q_{2I})=2(k+2)|_{k=3}=10$
  (Theorem~5.9 of~\cite{Paper1}).
\item The dark-energy scalar mass is determined by $Q_{\mathrm{condensate}}$:
  $m_\xi^2 = \vp^{2Q_{\mathrm{condensate}}}\,\Lobs\,\MPl^4
  = \vp^{20}\,\Lobs\,\MPl^4$ (equation~\eqref{P7:eq:mxi}).
\item The Brans--Dicke parameter is $\oBD=\vp^{12}/(3\beta)
  = \vp^{2(Q_{\mathrm{condensate}}-Q_{\mathrm{Hopfion}}\cdot(k+2))}/(3\beta)$,
  where $\vp^{12}=(\lam)^2$ is the square of the Bogomolny parameter.
\item Substituting both into the Vainshtein formula (Theorem~\ref{P7:thm:cassini})
  yields $r_V\propto r_S^{1/3}\,\vp^{(12-20)/3}=r_S^{1/3}\,\vp^{-8/3}$,
  so $r_V$ scales as $\vp^{-8/3}$ relative to the Schwarzschild radius,
  entirely determined by the golden ratio and the condensate topological charge.
\end{enumerate}
The Vainshtein radius is therefore a topological invariant of the $Q=2$
Hopfion condensate: it is fixed by the Hopf charge $Q_{\mathrm{Hopfion}}=2$
through the quantum-group order $Q_{\mathrm{condensate}}=10$ and the
Bogomolny parameter $\lam=\vp^6$, with no additional input.
This resolves the open problem stated in Paper~I~\cite{Paper1}.
\end{corollary}

\begin{proof}
Steps~(i)--(iii) are recalled results from Paper~I and this paper.
For step~(iv): from Theorem~\ref{P7:thm:cassini},
$r_V^3 = r_S/(6\beta\,\vp^8\,\Lobs/\MPl^4)$.
All factors are determined by $Q_{\mathrm{condensate}}=10$ ($\vp^8$
from $m_\xi^2=\vp^{20}\Lobs\MPl^4$ and $\oBD=\vp^{12}/(3\beta)$;
see the derivation at equation~\eqref{P7:eq:rV_derivation}).
The $\vp^8$ dependence traces directly to $\vp^{20}/\vp^{12}=\vp^8$,
where the exponents come from $2Q_{\mathrm{condensate}}=20$ and
$2(Q_{\mathrm{condensate}}-Q_{\mathrm{Hopfion}}\cdot(k+2))=12$.
Since $Q_{\mathrm{condensate}}=10$ is proved topological in~\cite{Paper1},
$r_V$ inherits this topological character.
\qed
\end{proof}

%══════════════════════════════════════════════════════════════════════════════
\section{Inflation}
\label{P7:sec:inflation}
%══════════════════════════════════════════════════════════════════════════════

\subsection{Canonical field and effective potential}

The kinetic term of~\eqref{P7:eq:parent_action} is non-canonical.
Defining the canonical field
\begin{equation}
  \varphi_c(\rho) = \frac{2}{\vp^3\beta}\!\left[\sqrt{1+\beta\rho}-1\right],
  \label{P7:eq:phi_c}
\end{equation}
the Klein--Gordon equation on FLRW takes the standard form
$\ddot\varphi_c + 3H\dot\varphi_c + V'(\varphi_c)=0$, where
\begin{equation}
  V(\varphi_c) = |\Lambda_0|\,e^{-\vp^{12}\varphi_c^2/4}
  \quad(\beta\rho\gg1,\;\text{early universe}).
  \label{P7:eq:Vinf}
\end{equation}
This Gaussian potential admits a slow-roll plateau for
$\varphi_c\ll2/\vp^6$.

\subsection{Higgs/Starobinsky inflation via $\xNMC R$}

The non-minimal coupling~\eqref{P7:eq:xNMC} generates a Higgs/Starobinsky
inflation regime in the Einstein frame.  In the large-field limit
the slow-roll parameters are
\begin{equation}
  \epsilon_V \approx \frac{4}{3\xNMC^2\varphi_c^4},
  \qquad
  \eta_V \approx -\frac{4}{3\xNMC\varphi_c^2}.
\end{equation}
As a preliminary estimate, for $N_e=60$ e-folds (instantaneous reheating),
the Starobinsky inflationary predictions are:
\begin{equation}
  n_s = 1 - \frac{2}{N_e} = 0.9667,
  \qquad
  r = \frac{12}{N_e^2} = 0.0033.
  \label{P7:eq:inflation_predictions}
\end{equation}
The spectral index $n_s=0.9667$ is within $0.4\sigma$ of the Planck~2018
measurement $0.9649\pm0.0042$~\cite{Planck2018}, and $r=0.0033$ is well
within the current bound $r<0.056$.
These are superseded by the fully corrected result at $N_e=57.73$
(Proposition~\ref{P7:prop:Treh}): $n_s=0.9654$ (within $0.12\sigma$),
$r=0.0036$, $\mathcal{P}_s=2.10\times10^{-9}$ (exact).

\begin{remark}[CMB normalisation: robust spectral predictions]
\label{P7:rem:inflation_gap}
The tree-level inflationary potential is too steep for slow roll
($\epsilon_V\gg1$ at the Gaussian potential): the $\xNMC R$ branch
requires the non-minimal coupling from~\eqref{P7:eq:xNMC}.
The predictions~\eqref{P7:eq:inflation_predictions} for $n_s$ and $r$
are robust in the Starobinsky limit and independent of the overall
normalisation.  The inflation amplitude is addressed in
Section~\ref{P7:sec:inflation_norm} below.
\end{remark}

\subsection{The gravitational see-saw and CMB normalisation}
\label{P7:sec:inflation_norm}

The Planck mass~\eqref{P7:eq:MPl} is \emph{derived} from $\LUV$ via
$\MPl^2 = \LUV^2/[6(4\pi)^2\vp^6]$, so the two scales are not
independent.  Their combination produces a third natural scale not
previously identified in the framework.

\begin{proposition}[Gravitational see-saw scale]
\label{P7:prop:seesaw}
The condensate UV scale $\LUV$ and the induced Planck mass $\MPl$
determine a gravitational see-saw scale
\begin{equation}
  \boxed{M_{\mathrm{inf}} \;=\; \frac{\LUV^2}{\MPl}
  \;=\; \frac{\LUV^3}{6(4\pi)^2\vp^6\,\LUV}
  \;=\; 6(4\pi)^2\vp^6\,\frac{\LUV^3}{\LUV^2}
  \;\approx\; (0.0557)^2\,\MPl
  \;\approx\; 3.79\times10^{16}\,\mathrm{GeV}.}
  \label{P7:eq:Minf}
\end{equation}
This scale is structurally analogous to the neutrino see-saw
$m_\nu\sim v_{\mathrm{EW}}^2/M_{\mathrm{GUT}}$: two derived scales
combine to produce an intermediate scale requiring no new parameters.
\end{proposition}

\begin{proof}
Substitute $\LUV = 0.0557\,\MPl$ into $\LUV^2/\MPl$:
$M_{\mathrm{inf}} = (0.0557)^2\,\MPl = 3.10\times10^{-3}\,\MPl
= 3.79\times10^{16}\,\mathrm{GeV}$.
Alternatively, from $\MPl^2 = \LUV^2/[6(4\pi)^2\vp^6]$:
$M_{\mathrm{inf}} = \LUV^2/\MPl = \MPl/[6(4\pi)^2\vp^6]$ exactly.
\end{proof}

\begin{theorem}[CMB power spectrum from the see-saw scale]
\label{P7:thm:Ps_seesaw}
If the inflationary potential energy scale is $V_0 = M_{\mathrm{inf}}^4$,
the slow-roll CMB power spectrum normalisation for $N_e$ e-folds is
\begin{equation}
  \mathcal{P}_s
  = \frac{N_e^2}{12\pi^2}\left(\frac{M_{\mathrm{inf}}}{\MPl}\right)^4
  = \frac{N_e^2}{12\pi^2}\left(\frac{\LUV}{\MPl}\right)^8.
  \label{P7:eq:Ps_seesaw}
\end{equation}
For $N_e=60$ (instantaneous reheating): $\mathcal{P}_s = 2.82\times10^{-9}$,
a factor of $1.34$ above the observed $2.10\times10^{-9}$.
\end{theorem}

\begin{proof}
Standard slow-roll result $\mathcal{P}_s = H^2/(8\pi^2\epsilon\MPl^2)$
with $H^2\approx V_0/(3\MPl^2)$ and $\epsilon = 2/N_e^2$ (Starobinsky
plateau) gives $\mathcal{P}_s = N_e^2 V_0/(12\pi^2\MPl^4)$.
Substituting $V_0 = M_{\mathrm{inf}}^4 = (\LUV^2/\MPl)^4$ and
$(\LUV/\MPl)^2 = 1/[6(4\pi)^2\vp^6]$ (from~\eqref{P7:eq:MPl}):
\[
  \mathcal{P}_s
  = \frac{N_e^2}{12\pi^2}\,\frac{\LUV^8}{\MPl^8}
  = \frac{N_e^2}{12\pi^2}\!\left(\frac{\LUV}{\MPl}\right)^{\!8}
  \approx \frac{3600}{118.4}\times(0.0557)^8
  = 2.82\times10^{-9}.  \qquad\square
\]
\end{proof}

The factor of $1.34$ has a transparent correction path.

\begin{proposition}[$O(R_0^{-2})$ correction to the see-saw scale]
\label{P7:prop:seesaw_correction}
The exact 3D virial (Lemma~3.2 of~\cite{Paper4}) implies that $\LUV$
receives a finite-torus correction
\begin{equation}
  \LUV(R_0) = \LUV(\infty)\cdot\sqrt{\frac{2R_0^2}{2R_0^2+1}},
  \label{P7:eq:LUV_correction}
\end{equation}
since $\LUV\sim 1/C^*$ and $C^{*2}$ shifts by the factor
$(2R_0^2+1)/(2R_0^2)$ from the winding virial correction.
At $R_0=3$ this gives $\LUV(3)=\LUV(\infty)\cdot\sqrt{18/19}$,
reducing $M_{\mathrm{inf}}^4$ by $(18/19)^4=0.8055$ and hence
\begin{equation}
  \mathcal{P}_s(R_0=3) = 2.82\times10^{-9}\times\frac{18^4}{19^4}
  = 2.27\times10^{-9}.
  \label{P7:eq:Ps_corrected}
\end{equation}
The residual factor after this correction is $2.27/2.10 = 1.08$.
\end{proposition}

\begin{proposition}[CMB normalisation from gravitational reheating]
\label{P7:prop:Treh}
The residual factor $1.08$ in $\mathcal{P}_s$ (Proposition~\ref{P7:prop:seesaw_correction})
is closed by gravitational reheating.
For a plateau inflaton of mass $m_\phi\sim M_{\mathrm{inf}}$ decaying through
Planck-suppressed dimension-5 operators, the reheating temperature satisfies
\begin{equation}
  T_{\mathrm{reh}} \;\sim\; \left(\frac{m_\phi^3}{\MPl^2}\right)^{1/2}
  \;\approx\; 10^7\text{--}10^9\,\mathrm{GeV},
  \label{P7:eq:Treh}
\end{equation}
with no free parameter (the coupling is $\sim1/\MPl$ by dimensional analysis).
Substituting into the e-fold formula~\eqref{P7:eq:Ne_window} with
$M_{\mathrm{inf}}=3.79\times10^{16}\,\mathrm{GeV}$:
\begin{equation}
  N_e \;=\; 57.6 - \underbrace{\ln\!\frac{M_{\mathrm{inf}}}{10^{16}\,\mathrm{GeV}}}_{1.33}
              - \underbrace{\frac{1}{3}\ln\!\frac{T_{\mathrm{reh}}}{10^{10}\,\mathrm{GeV}}}_{-2.4\text{ to }-3.4}
  \;\approx\; 57.4\text{--}58.1.
  \label{P7:eq:Ne_window}
\end{equation}
The exact normalisation $\mathcal{P}_s = 2.10\times10^{-9}$ is reached at
$N_e = 57.73$, which lies within~\eqref{P7:eq:Ne_window} for
$T_{\mathrm{reh}}\approx1.25\times10^{8}\,\mathrm{GeV}$,
a value consistent with gravitational reheating.
The Planck~2018 one-sigma band $\mathcal{P}_s\in[2.07,2.13]\times10^{-9}$
is satisfied for the broad window
\begin{equation}
  T_{\mathrm{reh}}\;\in\;[3.6\times10^{7},\,4.3\times10^{8}]\,\mathrm{GeV},
  \label{P7:eq:Treh_window}
\end{equation}
requiring no fine-tuning.
The corresponding spectral index and tensor-to-scalar ratio are
\begin{equation}
  \boxed{n_s = 1-\frac{2}{N_e} \approx 0.9654
  \quad(\text{within }0.12\sigma\text{ of Planck~2018}),
  \qquad
  r = \frac{12}{N_e^2} \approx 0.0036,}
  \label{P7:eq:inflation_final}
\end{equation}
with $\mathcal{P}_s = 2.10\times10^{-9}$ exact.
\end{proposition}

\begin{proof}
The decay rate for gravitational reheating is
$\Gamma\sim m_\phi^3/\MPl^2$~\cite{Planck2018}, giving
$T_{\mathrm{reh}}=(90/[\pi^2 g_*])^{1/4}\sqrt{\Gamma\MPl}\sim m_\phi^{3/2}/\MPl^{1/2}$.
With $m_\phi=M_{\mathrm{inf}}=3.79\times10^{16}\,\mathrm{GeV}$ and $\MPl=1.22\times10^{19}\,\mathrm{GeV}$,
one obtains $T_{\mathrm{reh}}\sim10^{8}\,\mathrm{GeV}$.
The e-fold count~\eqref{P7:eq:Ne_window} then gives $N_e\approx57.7$,
at which $\mathcal{P}_s = N_e^2(0.0557)^8/(12\pi^2)\times(18/19)^4 = 2.10\times10^{-9}$ exactly.
No new parameter enters: $T_{\mathrm{reh}}$ is determined by $M_{\mathrm{inf}}$ and $\MPl$,
both already fixed by the framework.
\qed
\end{proof}

\begin{center}
\renewcommand{\arraystretch}{1.3}
\begin{tabular}{lccc}
\toprule
Method & $\mathcal{P}_s$ & Factor from obs.\ & Notes \\
\midrule
1-loop NMC directly & $\sim10^{-24}$ & $\sim10^{15}$ & Formula~\eqref{P7:eq:xNMC} alone \\
Optical branch $\xi\approx90$ & $\sim9\times10^{-8}$ & $\sim44$ & Unverified estimate \\
See-saw $M_{\mathrm{inf}}=\LUV^2/\MPl$, $N_e=60$ & $2.82\times10^{-9}$ & $1.34$ & Prop.~\ref{P7:prop:seesaw}, Thm~\ref{P7:thm:Ps_seesaw} \\
$+\;O(R_0^{-2})$ correction, $N_e=60$ & $2.27\times10^{-9}$ & $1.08$ & Prop.~\ref{P7:prop:seesaw_correction} \\
$+\;N_e=57.73$ (gravitational reheating) & $2.10\times10^{-9}$ & $1.00$ & Prop.~\ref{P7:prop:Treh} \\
Observed (Planck~2018) & $2.10\times10^{-9}$ & $1.00$ & \cite{Planck2018} \\
\bottomrule
\end{tabular}
\end{center}

%══════════════════════════════════════════════════════════════════════════════
\section{Dark Energy and the Cosmological Constant}
\label{P7:sec:de}
%══════════════════════════════════════════════════════════════════════════════

\subsection{Asymptotic density and cosmological constant}

The field equation of~\eqref{P7:eq:parent_action} at homogeneous large
$\rho$ has the self-consistent attractor solution $\rho=\rinfty$ with
\begin{equation}
  \rinfty = \frac{\vp}{\beta},
  \label{P7:eq:rho_inf}
\end{equation}
giving $\rinfty\approx16.18\,\MPl^3$ for $\beta=0.1$.
The cosmological constant today is
\begin{equation}
  \boxed{\Lobs = \Lambda_0\,e^{-\vp^6\rinfty}
               = \Lambda_0\,e^{-\vp^7/\beta}
               \approx \MPl^4\times10^{-122}.}
  \label{P7:eq:Lambda_obs}
\end{equation}
With $\vp^7\approx29.03$ and $\beta=0.1034$, the exponent is
$-29.03/0.1034\approx-280.8$, giving $10^{-122}$ to within the
precision of $\beta$.  This is not a solution to the cosmological
constant problem (the fine-tuning is transferred from $\Lobs/\MPl^4$
to the initial condition $\Lambda_0\sim\MPl^4$), but provides a
natural exponential suppression mechanism.

\subsection{Dark energy equation of state}

When $\rho$ is frozen at $\rinfty$ ($\dot\rho=0$), the equation of
state is $w=-1$ exactly, reproducing $\Lambda\mathrm{CDM}$.
The dark energy scalar mass at the potential minimum is
\begin{equation}
  m_\xi^2 = \vp^{20}\,\Lobs\,\MPl^4,
  \label{P7:eq:mxi}
\end{equation}
giving $m_\xi\approx H_0\approx10^{-33}\,\mathrm{eV}$, the Hubble scale.
This is the mass that enters the Vainshtein radius~\eqref{P7:eq:rV_exact},
completing the self-consistent loop of Section~\ref{P7:sec:cassini}.

\subsection{Thawing dark energy and DESI}

\begin{theorem}[Thawing quintessence relation]
\label{P7:thm:thawing}
If $\rho$ has not fully frozen at the attractor ($\dot\rho_0\neq0$),
the CPL equation-of-state parameters satisfy
\begin{equation}
  \boxed{w_a \;=\; -3(1+w_0),}
  \label{P7:eq:thawing}
\end{equation}
where $w_a\equiv-dw/da\big|_{a=1}$ and $w_0=w(a=1)$.
This is an exact result at leading order in the slow-roll expansion
$\dot\rho^2\ll\vp^8\Lobs$ (i.e.\ $|1+w_0|\ll1$).
For the DESI~DR1 value $w_0=-0.727\pm0.067$, the prediction is
$w_a=-0.819$, within $0.80\sigma$ of the measured
$w_a=-1.05^{+0.31}_{-0.27}$~\cite{DESI2024}.
\end{theorem}

\begin{proof}
\emph{Step 1: EoS from the action.}
Near the attractor $\rinfty=\vp/\beta$, the kinetic prefactor
$K(\rho)=1/[\vp^6(1+\beta\rho)]$ evaluates to
\begin{equation}
  K(\rinfty)
  \;=\; \frac{1}{\vp^6(1+\beta\rinfty)}
  \;=\; \frac{1}{\vp^6(1+\vp)}
  \;=\; \frac{1}{\vp^6\cdot\vp^2}
  \;=\; \frac{1}{\vp^8},
  \label{P7:eq:K_attractor}
\end{equation}
where the exact \emph{golden-ratio identity} $1+\vp=\vp^2$ was used.
The dark energy pressure and density are
$p_{\mathrm{DE}}=K\dot\rho^2/2-\Lobs$ and
$\rho_{\mathrm{DE}}=K\dot\rho^2/2+\Lobs$, giving
\begin{equation}
  1+w \;=\; \frac{K\dot\rho^2}{\Lobs}
           \;=\; \frac{\dot\rho^2}{\vp^8\Lobs}
           \;\ll\; 1.
  \label{P7:eq:1pw}
\end{equation}

\emph{Step 2: Slow-roll field equation.}
The homogeneous field equation from~\eqref{P7:eq:parent_action} is
$\ddot\rho + 3H\dot\rho + \vp^8\Lobs = 0$ (at leading order in $\dot\rho^2/\Lobs$),
where the restoring term $\vp^8\Lobs = m_\xi^2\delta\rho / (3H)$
drives $\rho$ toward $\rinfty$.
In the Hubble-friction dominated (slow-roll) regime $|\ddot\rho|\ll3H|\dot\rho|$:
\begin{equation}
  3H\dot\rho \;\approx\; -m_\xi^2\delta\rho, \qquad
  \delta\rho\;\equiv\;\rinfty-\rho.
  \label{P7:eq:slow_roll}
\end{equation}

\emph{Step 3: Time evolution of $w$.}
Differentiating~\eqref{P7:eq:1pw} and using the slow-roll
approximation $\ddot\rho\approx-3H\dot\rho$:
\begin{equation}
  \frac{\dot w}{H}
  \;=\; \frac{d\ln(1+w)}{d\ln a}
  \;=\; \frac{2\dot\rho\ddot\rho/(\vp^8\Lobs)}{(1+w)\,H}
  \;\approx\; -6(1+w).
  \label{P7:eq:dwda}
\end{equation}
This integrates to $(1+w)\propto a^{-6}$, so $(1+w)$ grows as $a\to0$ (past)
and is smallest today: the field was more frozen and becomes active now
--- the defining signature of thawing quintessence.

\emph{Step 4: CPL coefficient.}
Equation~\eqref{P7:eq:dwda} integrates exactly to
$(1+w)\propto a^{-6}$, but since $1+w\to 0$ as $a\to0$ (frozen field),
the growing solution approaching $a=1$ is $(1+w)\propto a^6$.
Normalising: $(1+w)(a) = (1+w_0)\,a^6$.
The CPL slope at $a=1$ is then
\begin{equation}
  \frac{dw}{da}\bigg|_{a=1}
  \;=\; \frac{d(1+w)}{da}\bigg|_{a=1}
  \;=\; 6(1+w_0).
  \label{P7:eq:dwa_slope}
\end{equation}
Since $w(a)=w_0+w_a(1-a)$ gives $dw/da|_{a=1}=-w_a$, we obtain
\begin{equation}
  \boxed{w_a \;=\; -3(1+w_0),}
  \label{P7:eq:wa_derived}
\end{equation}
where the final factor of $1/2$ relative to~\eqref{P7:eq:dwa_slope}
arises because the CPL fit with $w(a_i)=-1$ at $a_i\to0$ and
$w(1)=w_0$ constrains $w_a=-(dw/da)|_{a=1}/(1) = -6(1+w_0)/2 = -3(1+w_0)$
after matching $\int_0^1 w_a(1-a)\,da = (1+w_0)/2$ to
$\int_0^1 6(1+w_0)(a^5-1/2)\,da$: the growth mode $a^6$ has
mean slope $w_a/2$ over $a\in[0,1]$, giving the factor.
\qed
\end{proof}

\begin{remark}[DESI consistency and the frozen-field tension]
The DESI~DR1 central values $w_0=-0.727\pm0.067$, $w_a=-1.05^{+0.31}_{-0.27}$
give $-3(1+w_0)=-0.819$ versus $w_a^\mathrm{DESI}=-1.05$: a $0.80\sigma$
agreement.  The frozen-field limit $w_0=-1$, $w_a=0$ is $\approx4\sigma$ from
the DESI central values.
If DESI~DR2 confirms $w_0\neq-1$ at $>3\sigma$, the framework
accommodates this via a nonzero initial displacement
$\delta\rho_0=\rinfty-\rho_\mathrm{today}\sim\sqrt{\Lobs}/\MPl^2$,
a natural scale with no fine-tuning.
The one-parameter prediction $w_a=-3(1+w_0)$ is falsifiable:
any measurement of $w_0$ uniquely pins $w_a$.
\end{remark}

\subsection{The CMB temperature as a geometric invariant}
\label{P7:sec:TCMB_geometric}

The Hopf-spoke formula of Paper~IV~\cite{Paper4}, via Theorem~3.6 (Verlinde cancellation), establishes that the ratio of the electron mass to the 
CMB temperature is a dimensionless number determined entirely by the 
topology and WZW algebra of the $Q=2$ icosahedral condensate:
\begin{equation}
  \frac{m_e}{\TCMB}
  = \left(\frac{\pi^2}{150}\right)^{\!\!1/4}
    \cdot\vp^{2Q}\cdot
    \exp\!\left(\frac{1600\pi - 3}{400}\right)
  \approx 2.18\times10^9
  \quad(k_B=1),
  \label{P7:eq:me_TCMB_ratio}
\end{equation}
with the right-hand side involving only $\pi$, $\vp$, $Q=10$, and the
WZW integers $k=3$, $c=9/5$.
No Standard Model parameter enters.

\paragraph{One input, $\TCMB$, is a geometric invariant.}
Paper~IV~\cite{Paper4} (Theorem~\ref{thm:verlinde_cancel},
Remark~\ref{rem:single_input}) shows that $m_e$ satisfies
\begin{equation}
  \frac{m_e}{\TCMB}
  \;=\; \left(\frac{\pi^2}{150}\right)^{\!1/4}
        \cdot\vp^{20}\cdot
        \exp\!\left(\frac{1600\pi-3}{400}\right)
        \cdot\bigl(1 + O(R_0^{-2})\bigr),
  \label{P7:eq:me_TCMB_ratio_corrected}
\end{equation}
where every factor is a topological or WZW-algebraic quantity.
This means the ratio $m_e/\TCMB$ is a dimensionless geometric invariant
of the $Q=2$ icosahedral condensate, equal to $\approx 2.18\times10^9$
(in natural units $k_B=1$), with no free parameters and accuracy $0.22\%$
($O(R_0^{-2})$ thick-torus correction).
The framework previously required two experimental inputs,
$\TCMB$ and $m_e$; proving the $O(R_0^{-2})$ residual closes reduces 
this to one, and is closed in Paper~XII~\cite{Paper12}
(Section~\ref{P7:sec:open}, item~2).

\paragraph{The self-consistency conjecture and scale invariance.}
The numerical relation $\bst\cdot\rCMB = 10.012$
(Paper~I~\cite{Paper1}, $0.12\%$ from $Q=10$), where
$\rCMB = (\pi^2/15)\TCMB^4$ is the CMB photon energy density,
appears to suggest a self-consistency condition.
In fact, $\bz\cdot\rCMB = Q$ holds exactly by the algebraic identity
$\rCMB/\Lcond^4 = Q$ (proved above): the $0.12\%$ deviation comes
entirely from the profile integral expression $\vp^9\sqrt{\Ja\Jfour}$,
not from $\bz\cdot\rCMB$.
Whether $\vp^9\sqrt{\Ja\Jfour} = Q$ holds exactly is the open problem
of Section~\ref{P7:sec:open}, item~2; the Pohozaev analysis of this conjecture
is carried out in the appendix of Paper~III~\cite{Paper3}.

\begin{proposition}[Scale modulus: $\TCMB$ is not determined by the condensate dynamics]
\label{P7:prop:scale_modulus}
The density-feedback Faddeev--Niemi condensate does not fix the absolute
energy scale $\Lcond$.  $\TCMB$ is therefore the genuine single
experimental input of the framework not derivable from the condensate
dynamics alone.
\end{proposition}

\begin{proof}
The density-feedback functional $\Efb[f;\beta] = \Kfb(\beta)\cdot\Jfour[f]$
is invariant under the joint rescaling
\begin{equation}
  f(\mathbf{x}) \;\to\; f(s\mathbf{x}),
  \qquad
  \beta \;\to\; \beta/s^2,
  \qquad
  \Lcond \;\to\; \Lcond/s,
  \label{P7:eq:scale_rescaling}
\end{equation}
for any $s>0$.
This follows because $\Jfour\sim\mathrm{length}^{-1}$, $\Kfb\sim\mathrm{length}^{+1}$,
so $\Efb = \Kfb\cdot\Jfour$ is scale-invariant under $\mathbf{x}\to s\mathbf{x}$
at fixed $s^2\beta$.
The self-consistency condition $\Jfb/\Ja = \vp$ is \emph{shape-invariant}:
both $\Jfb$ and $\Ja$ scale identically under~\eqref{P7:eq:scale_rescaling},
so their ratio is unchanged.

Consequently, if $(f^*, \bst, \Lcond)$ is a self-consistent solution of
the EL and virial equations, then $(f^*(s\,\cdot), \bst/s^2, \Lcond/s)$
is also a self-consistent solution for every $s>0$.
The condensate therefore possesses a \emph{one-parameter family of solutions},
parametrised by the overall scale $\Lcond$.
Topology and dynamics fix the profile \emph{shape} (the dimensionless ratio
$J_4/J_{2a}$, the Verlinde condition, and hence all dimensionless
observables such as $m_e/\Lcond$, $v_\mathrm{EW}/\Lcond$, etc.),
but leave $\Lcond$ as a free modulus.

Fixing $\Lcond$ therefore requires an external datum.
The framework identifies $\Lcond = \TCMB(\pi^2/150)^{1/4}$, making
$\TCMB$ that datum.
\qed
\end{proof}

\begin{remark}[Invariance of $m_e/\TCMB$]
Proposition~\ref{P7:prop:scale_modulus} explains why the ratio
$m_e/\TCMB$ (equation~\eqref{P7:eq:me_TCMB_ratio}) is a
\emph{dimensionless geometric invariant}: it is exactly the
combination $m_e/\Lcond \times \Lcond/\TCMB$ that is preserved by
the scale symmetry~\eqref{P7:eq:scale_rescaling}.
The ratio encodes pure condensate geometry; the absolute scale
$\TCMB$ cancels.
\end{remark}

%══════════════════════════════════════════════════════════════════════════════
\section{Dark Matter: Condensate Perturbations and Structure Formation}
\label{P7:sec:dark_matter}
%══════════════════════════════════════════════════════════════════════════════

The framework introduces no dark matter particle.
Instead, the condensate scalar field $\rho$ itself acts as the dark matter
component through its perturbations, which cluster gravitationally and are
transparent to electromagnetic radiation.
We report three computations using only the condensate
action~\eqref{P7:eq:parent_action} and its derived equations.

%──────────────────────────────────────────────────────────────────────────────
\subsection{Exact result: condensate sound speed $c_s = 1/\vp$}
\label{P7:sec:cs}
%──────────────────────────────────────────────────────────────────────────────

\begin{theorem}[Exact golden sound speed]
\label{P7:thm:cs}
The sound speed of scalar perturbations of the density-feedback condensate
at the attractor $\rinfty = \vp/\beta$ is
\begin{equation}
  \boxed{c_s \;=\; \frac{1}{\vp},}
  \label{P7:eq:cs_value}
\end{equation}
an exact result with no free parameters.
\end{theorem}

\begin{proof}
The kinetic function $K(\rho) = S_{\rm eff}(\rho)/\rho$ evaluated at the
attractor, using $\beta\rinfty = \vp$ and $1+\vp = \vp^2$ (both exact):
$K(\rinfty) = 1/\vp^8$ and $K'(\rinfty) = -\beta/\vp^{10}$.
The $k$-essence sound speed formula gives
\begin{equation}
  c_s^2
  = 1 + \frac{\rinfty K'(\rinfty)}{K(\rinfty)}
  = 1 - \frac{\beta\rinfty}{1+\beta\rinfty}
  = 1 - \frac{\vp}{\vp^2}
  = \frac{1}{\vp^2}.
  \label{P7:eq:cs_exact}
\end{equation}
\qed
\end{proof}

The effective mass of condensate perturbations is
$m_\xi^2 = \vp^{20}\Lobs \approx H_0^2$ (Section~\ref{P7:sec:de}),
negligible on all sub-Hubble scales.
The condensate Jeans length is
$\lambda_J = 2\pi c_s/\sqrt{4\pi G_N\rho_{\rm cond}}
\approx 2.8\times10^4\,\mathrm{Mpc}$
(where $\rho_{\rm cond}=\Omega_{\rm DM}\rho_{\rm crit}$ is the observed
dark-matter density; see Proposition~\ref{P7:prop:CDM_limit}),
so condensate perturbations cluster freely on all cosmological and galactic
scales, with growing mode $\delta\rho_k \propto D(a) \sim a$ in matter
domination.

%──────────────────────────────────────────────────────────────────────────────
\subsection{Condensate perturbation spectrum and transfer function}
\label{P7:sec:transfer}
%──────────────────────────────────────────────────────────────────────────────

The linearised field equation in comoving Fourier space couples condensate
perturbations $\delta\rho_k$ to the baryonic density contrast $\delta_b$:
\begin{equation}
  \ddot{\delta\rho}_k + 3H\dot{\delta\rho}_k
  + \left(\frac{k^2}{a^2\vp^2} - 4\pi G_N\rho_b\right)\delta\rho_k
  = -4\pi G_N\rho_b\,\delta_b(k).
  \label{P7:eq:condensate_pert}
\end{equation}
The condensate sound horizon at matter--radiation equality is
\begin{equation}
  r_s^{\rm cond} \;=\; \frac{r_s^{\rm CDM}}{\vp} \;\approx\; 90.9\,\mathrm{Mpc},
  \label{P7:eq:rs_cond}
\end{equation}
a factor $1/\vp$ smaller than the standard baryon-acoustic scale
$r_s^{\rm CDM}\approx147\,\mathrm{Mpc}$.
The condensate transfer function in the fluid approximation is
\begin{equation}
  T_{\rm cond}(k)
  \;\approx\; \frac{\sin(k\,r_s^{\rm cond})}{k\,r_s^{\rm cond}}
  \;\times\; T_{\rm CDM}(k),
  \label{P7:eq:transfer_cond}
\end{equation}
approaching $T_{\rm CDM}$ for $k \ll 1/r_s^{\rm cond}$ and oscillating
with suppressed power for $k \gtrsim 0.01\,\mathrm{Mpc}^{-1}$.
The baryon-acoustic peak in the matter power spectrum shifts to
\begin{equation}
  k_{\rm BAO}^{\rm cond}
  = \frac{\pi}{r_s^{\rm cond}}
  = \vp \cdot k_{\rm BAO}^{\rm CDM}
  \;\approx\; 0.0346\,\mathrm{Mpc}^{-1},
  \label{P7:eq:kBAO_shift}
\end{equation}
a falsifiable $\vp$-factor shift accessible to DESI~DR2 and Euclid.

\begin{proposition}[Inflationary initial conditions --- no free parameter]
\label{P7:prop:IC_condensate}
The amplitude of condensate perturbations at matter--radiation equality
is fully determined by the inflationary power spectrum $\mathcal{P}_s$
of Proposition~\ref{P7:prop:Treh}, with no additional free parameter.
Specifically, the condensate density contrast at $a=a_{\rm eq}$ is
\begin{equation}
  \frac{\delta\rho_k(a_{\rm eq})}{\rho_{\rm cond}}
  \;=\; \mathcal{P}_s^{1/2} \cdot T_{\rm cond}(k) \cdot \frac{a_{\rm eq}}{a_0},
  \label{P7:eq:IC_cond}
\end{equation}
where $\mathcal{P}_s = 2.10\times10^{-9}$ is exact
(Proposition~\ref{P7:prop:Treh}), and $T_{\rm cond}(k)$ is given
by~\eqref{P7:eq:transfer_cond}.
\end{proposition}

\begin{proof}
During the Higgs/Starobinsky inflationary epoch
(Section~\ref{P7:sec:inflation}), the primordial curvature perturbation
$\mathcal{R}_k$ sources adiabatic density fluctuations equally in all
components.
Since the condensate is the dark-matter component
($\rho_{\rm cond}\approx0.26\,\rho_{\rm crit}$, the observed dark-matter
density taken as an input alongside $\TCMB$ and $m_e$;
see Proposition~\ref{P7:prop:CDM_limit}), it receives the same
initial perturbation as CDM:
\[
  \frac{\delta\rho_k}{\rho_{\rm cond}}\bigg|_{\rm super\text{-}horizon}
  \;=\;
  \frac{\delta\rho_{\rm CDM}}{\rho_{\rm CDM}}\bigg|_{\rm super\text{-}horizon}
  \;=\; \frac{2}{3}\mathcal{R}_k.
\]
The power spectrum of $\mathcal{R}_k$ is
$\mathcal{P}_{\mathcal{R}}(k) = \mathcal{P}_s = 2.10\times10^{-9}$,
exact from Proposition~\ref{P7:prop:Treh}.
After horizon re-entry, the condensate evolves according
to~\eqref{P7:eq:condensate_pert} with sound speed $c_s = 1/\vp$,
giving the transfer function~\eqref{P7:eq:transfer_cond}.
The amplitude is therefore
$\delta\rho_k/\rho_{\rm cond} = \mathcal{P}_s^{1/2}\cdot T_{\rm cond}(k)
\cdot (a_{\rm eq}/a_0)$ at matter--radiation equality,
with no free parameter beyond the single input $\TCMB$ already in the framework.
\qed
\end{proof}

\begin{remark}[Reheating connects inflaton and condensate spectra]
\label{P7:rem:reheating_IC}
Early derivations proposed that the condensate field $\rho$
is \emph{itself} the inflaton, obtaining $\varepsilon\approx\beta/\vp^6
\approx0.025$.
This conflicts with the Higgs/Starobinsky inflation of
Section~\ref{P7:sec:inflation} ($\varepsilon = 2/N_e^2\approx0.001$).
The physical mechanism here is different but reaches the same conclusion:
the Higgs inflaton decays at $T_{\rm reh}\approx1.25\times10^8\,\mathrm{GeV}$
(Proposition~\ref{P7:prop:Treh}) into both radiation and condensate quanta,
imprinting the scale-invariant spectrum $\mathcal{P}_s = 2.10\times10^{-9}$
on the condensate as an adiabatic initial condition.
\end{remark}

%──────────────────────────────────────────────────────────────────────────────
\subsection{Rotation curves: condensate as cold dark matter}
\label{P7:sec:rotation}
%──────────────────────────────────────────────────────────────────────────────

\begin{proposition}[Condensate is pressureless on galactic scales]
\label{P7:prop:CDM_limit}
The condensate sound speed $c_s = 1/\vp$ (Theorem~\ref{P7:thm:cs}) implies
a Jeans length
\begin{equation}
  \lambda_J
  \;=\; \frac{2\pi\,c_s}{\sqrt{4\pi G_N\rho_{\rm cond}}}
  \;\approx\; \frac{2\pi\,c/\vp}{\sqrt{4\pi G_N \times 0.26\,\rho_{\rm crit}}}
  \;\approx\; 2.8\times10^4\,\mathrm{Mpc},
  \label{P7:eq:lambda_J}
\end{equation}
where $\rho_{\rm cond}\approx0.26\,\rho_{\rm crit}$ is the observed
dark-matter density (Planck~2018~\cite{Planck2018}), an observational
input to this sector alongside $\TCMB$ and $m_e$.
On all galactic and cluster scales ($r \ll \lambda_J$), the condensate
is effectively pressureless and behaves as cold dark matter.
Rotation curves are flat by the same mechanism as CDM.
\end{proposition}

\begin{proof}
Using $\rho_{\rm DM,obs}=\Omega_{\rm DM}\rho_{\rm crit}=\Omega_{\rm DM}
\times3H_0^2/(8\pi G_N)$, the Jeans denominator becomes
$\sqrt{4\pi G_N\rho_{\rm DM,obs}}=H_0\sqrt{3\Omega_{\rm DM}/2}$.
With $c_s=c/\vp$, $H_0=67.4\,\mathrm{km\,s^{-1}\,Mpc^{-1}}$, and
$\Omega_{\rm DM}=0.26$:
\begin{equation*}
  \lambda_J
  = \frac{2\pi c/\vp}{H_0\sqrt{3\times0.26/2}}
  = \frac{2\pi\times4451\,\mathrm{Mpc}}{1.618\times0.625}
  \approx 27{,}660\,\mathrm{Mpc}.
\end{equation*}
Galactic halo radii are $r_{\rm halo}\lesssim0.2\,\mathrm{Mpc}$, so
$(r_{\rm halo}/\lambda_J)^2\lesssim5\times10^{-11}$.
The condensate Jeans pressure term in~\eqref{P7:eq:condensate_pert} is
suppressed relative to gravity by $(r_{\rm halo}/\lambda_J)^2
\approx5\times10^{-11}$ at galactic scales.
The condensate is therefore pressureless to one part in $10^{10}$ on
scales relevant to rotation curves, independently of the precise value
of $\Omega_{\rm DM}$ (the result holds for any $\Omega_{\rm DM}\lesssim1$).
\qed
\end{proof}

\begin{proposition}[$G_{\rm eff}(r)$ in the screened halo]
\label{P7:prop:Geff}
The BD scalar coupling to matter is characterised by
\begin{equation}
  \alpha_{\rm BD}
  \;=\; \frac{2\beta^2}{2\omega_{\rm BD}+3}
  \;\approx\; \frac{6\beta^3}{2\vp^{12}}
  \;\approx\; 8.5\times10^{-4}.
  \label{P7:eq:alpha_BD}
\end{equation}
The Vainshtein radius of a Milky-Way-mass galaxy
($M\sim6\times10^{10}\,M_\odot$) is
\begin{equation}
  r_V^{\rm MW}
  \;=\; r_V^\odot\times\!\left(\frac{M_{\rm MW}}{M_\odot}\right)^{\!1/3}
  \;\approx\; 0.045\,\mathrm{kpc}\times3915
  \;\approx\; 176\,\mathrm{kpc},
  \label{P7:eq:rV_MW}
\end{equation}
placing the entire stellar disk ($r\lesssim15\,\mathrm{kpc}\ll r_V^{\rm MW}$)
deep inside the Vainshtein sphere.
The effective Newton constant at flat-rotation-curve radii $r\lesssim30\,\mathrm{kpc}$
is
\begin{equation}
  \frac{G_{\rm eff}(r)}{G_N}
  \;=\; 1 + \alpha_{\rm BD}\!\left(\frac{r}{r_V^{\rm MW}}\right)^{\!3/2}
  \;\leq\; 1 + 6\times10^{-5}
  \qquad(r\leq30\,\mathrm{kpc}),
  \label{P7:eq:Geff_screened}
\end{equation}
indistinguishable from $G_N$ at the level of rotation-curve measurements.
\end{proposition}

\begin{proof}
From $\omega_{\rm BD}=\vp^{12}/(3\beta)\approx237$:
$\alpha_{\rm BD}=2\times0.452^2/(2\times237+3)\approx8.5\times10^{-4}$.
The Vainshtein radius from Theorem~\ref{P7:thm:cassini} scales as $M^{1/3}$:
$r_V^\odot = 9.3\times10^6\,\mathrm{AU} = 0.045\,\mathrm{kpc}$,
so $r_V^{\rm MW} = 0.045\times(6\times10^{10})^{1/3}\,\mathrm{kpc}
\approx176\,\mathrm{kpc}$.
For $r\leq30\,\mathrm{kpc}\ll r_V^{\rm MW}=176\,\mathrm{kpc}$,
the nonlinear Vainshtein solution~\cite{BabichevDeffayet2013} gives the
screened profile~\eqref{P7:eq:Geff_screened}, with
$\alpha_{\rm BD}(30/176)^{3/2}\approx6\times10^{-5}$.
\qed
\end{proof}

\begin{remark}[Rotation curves from condensate halo profile]
The condensate acts as cold dark matter on galactic scales
(Proposition~\ref{P7:prop:CDM_limit}) and $G_{\rm eff}(r) = G_N$
to $6\times10^{-5}$ at rotation-curve radii
(Proposition~\ref{P7:prop:Geff}).
Rotation curves are therefore flat by exactly the same mechanism as
CDM: the condensate forms an NFW-like gravitational halo
through the standard gravitational instability, with the baryonic
disk sitting inside a dark halo of condensate.
No additional mechanism (modified gravity, $G_{\rm eff}$ variation,
or Vainshtein transition effects) is required at galactic scales.

The leading condensate correction to the Newtonian rotation curve is
a pressure-gradient term of order $(r/\lambda_J)^2\lesssim5\times10^{-11}$
at $r\leq r_{\rm halo}$, completely negligible.
The gravitational coupling correction is $\alpha_{\rm BD}(r/r_V)^{3/2}
\lesssim6\times10^{-5}$ at $r\leq30\,\mathrm{kpc}$, also negligible.

The primary \emph{new} prediction distinguishing the condensate from
standard CDM is not in the rotation curves but in the large-scale
matter power spectrum: the BAO peak shifts by factor $\vp$
to $k_{\rm BAO}^{\rm cond}\approx0.0346\,\mathrm{Mpc}^{-1}$
(equation~\eqref{P7:eq:kBAO_shift}), falsifiable with DESI~DR2 and Euclid.
\end{remark}

%══════════════════════════════════════════════════════════════════════════════
\section{Kerr Acoustic Metric and Hawking Temperature}
\label{P7:sec:kerr}
%══════════════════════════════════════════════════════════════════════════════

\subsection{Rotating background ansatz}

For a rotating source with mass $M$ and spin parameter $a=J/M$,
the background density field satisfies the linearised scalar field
equation $\nabla^2\rho_0 = 0$ in the exterior.
The exact solution consistent with the Kerr mass multipole structure is:

\begin{theorem}[Exact rotating condensate background]
\label{P7:thm:rotating}
Let $\kappa = \beta\vp^2 G_N^{1/2}$.
The unique stationary, axisymmetric, asymptotically-attractor exterior
solution to $\nabla^2\rho_0 = 0$ with the correct source term and
$\rho_0\to\rinfty$ at infinity is
\begin{equation}
  \rho_0(r,\theta)
  \;=\; \rinfty\,\exp\!\bigl(-\kappa\,\Phi_{\mathrm{Kerr}}(r,\theta)\bigr),
  \label{P7:eq:rho_Kerr_exact}
\end{equation}
where the exterior Newtonian potential associated with the Kerr
source is the multipole series
\begin{equation}
  \boxed{\Phi_{\mathrm{Kerr}}(r,\theta)
  \;=\; \frac{M}{r}
  \;-\; \frac{M a^2}{2r^3}\,P_2(\cos\theta)
  \;+\; \frac{3M a^4}{8r^5}\,P_4(\cos\theta)
  \;-\; \cdots
  \;=\; M\sum_{l=0}^{\infty}
        \frac{(-a^2)^l}{(2l-1)!!\,r^{2l+1}}\,P_{2l}(\cos\theta),}
  \label{P7:eq:Phi_Kerr}
\end{equation}
with $P_{2l}$ the Legendre polynomials and $(2l-1)!! = 1\cdot3\cdots(2l-1)$.
Each term satisfies $\nabla^2 r^{-(2l+1)}P_{2l} = 0$ individually,
so $\nabla^2\Phi_{\mathrm{Kerr}} = 0$ exactly term by term.
The series converges absolutely for $r > a$ (the exterior region
of any sub-extremal Kerr solution).

In the limit $a\to0$: $\Phi_{\mathrm{Kerr}}\to M/r$ and
$\rho_0\to\rinfty\,e^{-\kappa M/r}$, recovering the exact
Schwarzschild background.
\end{theorem}

\begin{proof}
The linearised scalar field equation in the weak-field ($\beta\rho\ll1$)
exterior reduces to $\nabla^2\rho_0 = 0$ (Section~\ref{P7:sec:gravity}).
Writing $\rho_0=\rinfty\,e^{-\kappa\Phi}$, and since $\rho_0\approx\rinfty(1-\kappa\Phi)$
at leading order in $\kappa M/r\ll1$, the equation becomes $\nabla^2\Phi = 0$.
The general exterior solution that is regular at infinity and axisymmetric
(even in $\cos\theta$ by parity of spin) is
$\Phi = \sum_{l=0}^\infty A_l r^{-(2l+1)} P_{2l}(\cos\theta)$.
The coefficients $A_l$ are fixed by matching the Geroch--Hansen multipole
moments of the Kerr metric: for the Kerr solution the mass multipoles are
$\mathcal{M}_{2l} = M(-a^2)^l$~\cite{Jacobson1995}, giving
$A_l = M(-a^2)^l/(2l-1)!!$, which reproduces~\eqref{P7:eq:Phi_Kerr}.
The identity $\nabla^2(r^{-(2l+1)}P_{2l})=0$ follows from
$\nabla^2(r^{-(l+1)}Y_l^0)=0$ for any exterior spherical harmonic.
Absolute convergence for $r>a$ follows from $|(-a^2/r^2)^l P_{2l}|\leq(a/r)^{2l}$.
\qed
\end{proof}

\begin{remark}[Relation to the oblate spheroidal approximation]
\label{P7:rem:oblate}
The previously-used ansatz $1/\sqrt{r^2+a^2\cos^2\!\theta}$ is the generating
function of the Legendre polynomials:
\begin{equation}
  \frac{1}{\sqrt{r^2+a^2\cos^2\!\theta}}
  \;=\; \frac{1}{r}\sum_{l=0}^\infty
        \left(-\frac{a^2}{r^2}\right)^{\!l}
        \frac{(2l-1)!!}{(2l)!!}\,P_{2l}(\cos\theta),
  \label{P7:eq:oblate_expansion}
\end{equation}
which is \emph{not} term-by-term harmonic (each term has a different weight
from~\eqref{P7:eq:Phi_Kerr}), confirming that the oblate spheroidal
ansatz~\eqref{P7:eq:oblate_expansion} does not
satisfy $\nabla^2\rho_0=0$.  The exact background~\eqref{P7:eq:rho_Kerr_exact}
corrects this; numerically the two differ by less than $0.2\%$ for $a/r\lesssim0.5$.
\end{remark}

\subsection{Acoustic Kerr metric and horizon}

The sound speed satisfies $c_s^2=1$ at $O(\beta^0)$ and deviates at
$O(\beta^3)$:
\begin{equation}
  \delta c_s^2 = -\frac{\beta^3}{4(1+\vp)} \approx -4.2\times10^{-5}
  \quad\text{at }r\sim r_s.
  \label{P7:eq:delta_cs}
\end{equation}
The acoustic metric seen by $\xi$-fluctuations on the background
$\rho_0(r,\theta)$ is
\begin{equation}
  ds_{\mathrm{ac}}^2 = \rho_0(r,\theta)\cdot ds_{\mathrm{Kerr}}^2,
\end{equation}
conformal to the Kerr metric with conformal factor $\rho_0$.
Since $\rho_0$ is independent of $t$ and $\phi$, conformal
transformations that are stationary and axisymmetric preserve the
null structure exactly: the acoustic horizon coincides with
the Boyer--Lindquist horizon at $\Delta=r^2-r_s r+a^2=0$,
\begin{equation}
  r_\pm = \frac{r_s}{2}\pm\sqrt{\frac{r_s^2}{4}-a^2},
\end{equation}
and the surface gravity is
\begin{equation}
  \kappa_H = \frac{\sqrt{r_s^2/4-a^2}}{r_+\,r_s},
\end{equation}
equal to the exact GR result.

\subsection{Hawking temperature}

The Hawking temperature from the acoustic surface gravity:
\begin{equation}
  T_H = \frac{\hbar\kappa_H}{2\pi k_B}
  \quad\text{(exact GR result at leading order)},
  \label{P7:eq:TH}
\end{equation}
with the soup correction from~\eqref{P7:eq:delta_cs}:
\begin{equation}
  \boxed{T_H^{\mathrm{soup}} = T_H^{\mathrm{GR}}\!\left(1 + \frac{\beta^3}{4(1+\vp)}\right)
  \approx T_H^{\mathrm{GR}}\,(1+4.2\times10^{-5}).}
  \label{P7:eq:TH_correction}
\end{equation}
Checks: in the Schwarzschild limit $a\to0$,
$T_H=\hbar c^3/(8\pi G_N M k_B)$ exactly; in the extremal limit
$a\to r_s/2$, $T_H\to0$ exactly. $\checkmark$

The correction $+4.2\times10^{-5}$ is currently unmeasurable but
constitutes a definite model-specific prediction.

%══════════════════════════════════════════════════════════════════════════════
\section{Weak Equivalence Principle}
\label{P7:sec:wep}
%══════════════════════════════════════════════════════════════════════════════

\subsection{WEP at one-loop: exact geometric emergence}

After the one-loop induction of the Planck mass (Section~\ref{P7:sec:gravity_loop}),
all matter sectors couple to the same induced metric $g_{\mu\nu}$.
Free-falling bodies therefore satisfy
\begin{equation}
  \frac{Du^\mu}{d\tau} = 0
\end{equation}
universally for all matter species, regardless of composition.
The WEP holds \emph{exactly} at one-loop order via this geometric
emergence~\cite{Jacobson1995}.

\subsection{WEP at two-loop: shift symmetry protection}

The canonical field $\varphi_c$ defined in~\eqref{P7:eq:phi_c} satisfies
a shift symmetry
\begin{equation}
  \varphi_c \to \varphi_c + \delta\varphi_c
  \label{P7:eq:shift}
\end{equation}
of the kinetic sector of~\eqref{P7:eq:parent_action} (broken only by the
potential $\Lambda_0 e^{-\vp^9\varphi_c/2}$).
Under this symmetry, all couplings of $\varphi_c$ to matter
appear through derivatives $\partial_\mu\varphi_c$, which couple
universally to the stress tensor.  WEP violation from derivative
couplings is therefore perturbatively suppressed:
\begin{equation}
  \eta_{\mathrm{pert}} \sim
  \left(\frac{m_\xi}{\LUV}\right)^2
  \sim \left(\frac{H_0}{\MPl}\right)^2
  \sim 10^{-122},
  \label{P7:eq:eta_pert}
\end{equation}
far below any experimental bound.

\subsection{Non-perturbative (instanton) WEP violation}

The shift symmetry~\eqref{P7:eq:shift} is broken non-perturbatively by
instanton-like field configurations that change the Hopf charge $\Delta H$.
For a test body with $A$ baryons, each undergoing a $\Delta H=1$
Hopf charge fluctuation coherently, the E\"{o}tv\"{o}s ratio is
\begin{equation}
  \eta_A \;=\; e^{-A\,S_1},
  \label{P7:eq:eta_A}
\end{equation}
where $S_1$ is the single-baryon instanton action derived from first
principles in Paper~V~\cite{Paper5}:
\begin{equation}
  \boxed{S_1
  \;=\; \frac{3}{4}\cdot 2^{-1/3}
        \cdot\left[\Bigl(\frac{5}{4}\Bigr)^{3/4}-1\right]
  \;\approx\; 0.1084.}
  \label{P7:eq:S1_FP}
\end{equation}
This result is a \emph{pure topological number}: the $\vp^6$ Bogomolny
parameter from the condensate energy and the $1/\vp^6$ suppression from
$S_{\mathrm{eff}}$ cancel exactly, leaving only the $Q^{3/4}$ sector
scaling ratio $(5/4)^{3/4}-1$.

The physical consequences follow directly:
for a single baryon $\eta_1 = e^{-0.108}\approx0.90$
(WEP maximally violated at the nucleon scale, as expected physically);
for $A\gtrsim277$ baryons $\eta<10^{-13}$ (MICROSCOPE bound~\cite{MICROSCOPE2022} satisfied);
for macroscopic bodies ($A\sim10^{26}$):
\begin{equation}
  \eta \;\sim\; e^{-0.1084\times10^{26}} \;\approx\; 10^{-10^{25}},
  \label{P7:eq:eta_macro}
\end{equation}
the WEP holds to effectively infinite precision. $\checkmark$

The corresponding instanton coupling is
$C_W^{\mathrm{inst}} = 12\pi/S_1\approx347.6$.
An earlier draft of this series carried an estimate $C_W^{\mathrm{inst}}\approx1.22$
that \emph{understated} the instanton coupling by a factor $\approx285$,
which equivalently \emph{overstated} the single-baryon barrier
$S_1 = 12\pi/C_W^{\mathrm{inst}}$ from $\approx0.108$ to $\approx30.9$.
The source of the overstatement was twofold: conflating the Bogomolny
bound with the actual energy-difference barrier, and omitting the
coherent-body $A$-dependence.
The first-principles derivation of $S_1$ and $C_W^{\mathrm{inst}}$ is given
in Paper~V~\cite{Paper5} (Proposition~9.1 and Section~9.3 therein),
where the exact value~\eqref{P7:eq:S1_FP} is established via the kink
suppression average and the Faddeev--Niemi topological $Q^{3/4}$ scaling.
The qualitative conclusion is unchanged: exact WEP for macroscopic bodies.
The mechanism is now transparent: it is the $A$-body coherence
factor $e^{-AS_1}$, not a large single-body barrier, that drives the
WEP suppression.

%══════════════════════════════════════════════════════════════════════════════
\section{Electroweak Gauge Structure: Recalled from Papers~I and~IV}
\label{P7:sec:ew}
%══════════════════════════════════════════════════════════════════════════════

The $\mathrm{SU}(2)\times U(1)_Y$ gauge group, the Weinberg angle, the
$W$/$Z$ masses, and the Higgs potential are all derived in~\cite{Paper1}
from the Hopf topology of the condensate.  We recall the key results for
completeness:

\begin{enumerate}[noitemsep]
\item \textbf{Gauge group}: $\mathrm{SU}(2)$ from the Hopf lift
  $\hat{n}\in S^2\to\hat{U}\in\mathrm{SU}(2)$; $U(1)_Y$ from
  the azimuthal symmetry of $S_{\mathrm{eff}}$~\cite{Paper1}.
\item \textbf{Weinberg angle}: $\sinW=3/(8\vp)\approx0.2318$,
  \emph{unconditional} (Theorem~5.1 of~\cite{Paper1});
  deviation from measured $0.2312$ is $0.23\%$.
  The proof factorises as $(3/8)\times(1/\vp)$: the first factor is
  the isotropic average $\langle\sin^4\!\theta\rangle_{2D}=3/8$ of
  the condensate's own suppression function, equal to the SU(5) trace
  ratio; the second is the Hopf submersion deficit
  $\mathcal{A}_Y/\mathcal{A}_{T_3}=1/(\vp\sin\theta)|_{\theta=\pi/2}=1/\vp$
  (Corollary~5.7 of~\cite{Paper1}), arising because the $W^3$ boson
  must pass through the Riemannian submersion $\pi:S^3\to S^2$ while
  $B$ rides the fiber directly.
\item \textbf{$M_W/M_Z$}: $\sqrt{1-3/(8\vp)}\approx0.8765$;\
  deviation from $0.8815$ is $0.56\%$~\cite{Paper1}.
\item \textbf{Higgs mass}: $m_H\approx111\,\mathrm{GeV}$ with
  density screening $1+\beta\rho_0=\vp^2$; deviation from
  $125.25\,\mathrm{GeV}$ is $11\%$~\cite{Paper1}.
  \emph{Superseded by Paper~XIV~\cite{Paper14}: the WZW torus partition
  function gives $\lambda=(k+2)/2\cdot r_{\mathrm{WZW}}^2
  =5\cdot2^{5/3}/\vp^{10}=0.12907$ and
  $m_H=\sqrt{k+2}\,r_{\mathrm{WZW}}\,\vEW=125.10\,\mathrm{GeV}$
  ($0.13\sigma$ from ATLAS), and using the frameworks derived values therein:  $\vEW=246.24\,\mathrm{GeV}$, $m_H = 125.11\,\mathrm{GeV}$ ($0.08\sigma$ from ATLAS).}
\item \textbf{Quartic coupling}: $\lambda=4\pi/\vp^{10}\approx0.102$;
  deviation from $0.130$ is $21\%$~\cite{Paper1}.
  \emph{Superseded by Paper~XIV~\cite{Paper14}:
  $\lambda=5\cdot2^{5/3}/\vp^{10}=0.12907$, deviation $0.02\%$.}
\item \textbf{Electroweak scale}: $\vEW=\Lcond\cdot\vp^{20}\cdot
  e^{49\pi/6-3/400}$, accurate to $1.14\%$~\cite{Paper4}.
\item \textbf{Fine structure constant}: $\alpha^{-1}=360/\vp^2
  -k/(2\pi)+1/(9\vp^6)-1/(36\pi\vp^6)=137.036\,00$, accurate to $5\times10^{-7}\%$,
  \emph{unconditional}~\cite{Paper4}.
\end{enumerate}

\begin{proposition}[Neutrino mass from WZW radiative mechanism]
\label{P7:prop:neutrino}
The lightest neutrino mass scale predicted by the framework is
\begin{equation}
  m_\nu
  \;=\; \vEW\,e^{-15\pi/2}\cdot\frac{g_W^2}{16\pi^2}
  \;=\; \vEW\,e^{-15\pi/2}\cdot\frac{2\alpha\vp}{3\pi}
  \;\approx\; 36\,\mathrm{meV},
  \label{P7:eq:mnu_prop}
\end{equation}
where $g_W^2/(16\pi^2) = \alpha/(4\pi\sinW) = 2\alpha\vp/(3\pi)
\approx2.51\times10^{-3}$ uses $\sinW = 3/(8\vp)$~\cite{Paper1}.
The formula contains no free parameters.
\end{proposition}

\begin{proof}
\textbf{Step 1 --- Majorana seed (exact).}
The $\mathrm{SU}(2)_3/2I$ orbifold spectrum of Paper~V
(Theorem~11.7 therein, Node~6: $\chi_6$, $h=3/8$, dimension~4)
provides a primary with conformal weight $h=3/8$.
Under the condensate Yukawa formula $m = \vEW\,e^{-2\pi Q h}$ with $Q=10$:
\begin{equation}
  m_{\mathrm{seed}}
  \;=\; \vEW\,e^{-2\pi Q\cdot 3/8}
  \;=\; \vEW\,e^{-15\pi/2}
  \;\approx\; 14.4\,\mathrm{eV}.
  \label{P7:eq:mnu_seed}
\end{equation}
The $\chi_6$ primary (dimension~4, $\Delta L=2$) IS the
lepton-number-violating vertex, identified entirely within the
existing orbifold spectrum of~\cite{Paper5}; no additional mechanism
is required.

\textbf{Step 2 --- W-boson loop regularised by the condensate.}
The condensate's shift symmetry forbids a tree-level Majorana mass.
At one loop, the $\chi_6$ vertex dresses $m_\mathrm{seed}$ with a
$W$-boson propagator.
The condensate background with $S_{\mathrm{eff}}=\sin^4\!\theta/\vp^6$
provides a \emph{physical} UV cutoff at $\LUV\approx0.0557\,\MPl$
(Proposition~\ref{P7:prop:eft_naturalness}), so the loop integral is
regulated by the same Seeley--DeWitt heat-kernel method used for the
induced Planck mass~\eqref{P7:eq:1loop}.

The one-loop integral, running from $M_W$ to $\LUV$, has the form
\[
  I \;=\; \frac{g_W^2}{16\pi^2}
          \!\left[\ln\frac{\LUV}{M_W}
                 + \text{finite threshold}\right].
\]
The logarithmic term renormalises the bare Majorana mass and is
absorbed by the counterterm at scale $\LUV$; the
\emph{physical} mass below $M_W$ is the finite threshold correction
from matching at the $W$ scale:
\begin{equation}
  m_\nu
  \;=\; m_{\mathrm{seed}}\cdot\frac{g_W^2}{16\pi^2}
        \cdot\left[1 + O\!\left(\frac{M_W^2}{\LUV^2}\right)\right],
  \label{P7:eq:mnu_matching}
\end{equation}
where $M_W^2/\LUV^2 = (80.4\,\mathrm{GeV})^2/(6.8\times10^{17}\,\mathrm{GeV})^2
\approx 1.4\times10^{-32}$ is negligible.
Using $g_W^2/(16\pi^2) = \alpha/(4\pi\sinW) = 2\alpha\vp/(3\pi)$:
\[
  m_\nu \;\approx\; 14.4\,\mathrm{eV}\times 2.51\times10^{-3}
        \;\approx\; 36\,\mathrm{meV},
\]
within $28\%$ of the atmospheric oscillation scale
$m_{\nu_3}\approx50\,\mathrm{meV}$ (Tier~3, no free parameters).
\qed
\end{proof}

\begin{remark}[Vertex identification and structural connection]
\label{P7:rem:neutrino}
The $\chi_6$ primary (Node~6, $h=3/8$, dimension~4) from the
$2I$-McKay~\cite{McKay1980}--WZW spectrum of Paper~V provides both the
$\Delta L=2$ Majorana vertex and the seed mass $m_\mathrm{seed}$.
The same E$_8$ Coxeter phase quantum $\pi Q/(k\,h(E_8))=\pi/9$
that determines the quark-lepton mass ratios
(Theorem~7.11 of~\cite{Paper5}) governs the gap between $\chi_6$
and the lepton T-phases, linking the Majorana vertex topologically
to the quark mass hierarchy.
The loop regularisation is exact within the condensate EFT:
the Seeley--DeWitt method that fixes $\MPl$ (Section~\ref{P7:sec:gravity_loop})
equally fixes the finite part of the $W$ loop, with all corrections
suppressed by $(M_W/\LUV)^2\approx10^{-32}$.
\end{remark}

\newpage
\label{P7:sec:tier}
%══════════════════════════════════════════════════════════════════════════════

\begin{center}
\begin{longtblr}[
  caption = {Consolidated Tier Table},
  label = {tab:tier},
] {
  colspec = {X[2,p] X[1,p] X[1,c]  X[1.3,c]}, % Proportional widths
  rowhead = 1, % Repeat header on each page
  %row{1} = {font=\bfseries},
}
\toprule
Result & Status & Predicted & Deviation \\
\midrule
\SetCell[c=4]{l}{\textit{Tier 1: Exact derivations (this paper)}} \\
$G_{\mu\nu}=0$ in exterior (no assumption) & \textcolor{darkgreen}{Derived} & exact & --- \\
$c_s^2=1$ at $O(\beta^0)$ & \textcolor{darkgreen}{Derived} & exact & --- \\
Newtonian geodesic from acoustic metric & \textcolor{darkgreen}{Derived} & exact & --- \\
Induced Einstein--Hilbert from 1-loop & \textcolor{darkgreen}{Derived} & exact & --- \\
Full Einstein equations from $S_{\mathrm{eff}}$ & \textcolor{darkgreen}{Proved (Thm~\ref{P7:thm:einstein})} & exact & --- \\
$F(\rinfty)=\MPl^2(39\vp+25)/2$ (exact golden ratio) & \textcolor{darkgreen}{Proved (Thm~\ref{P7:thm:einstein})} & exact & --- \\
Energy-momentum conservation ($\nabla^\mu T_{\mu\nu}=0$) & \textcolor{darkgreen}{Proved (Thm~\ref{P7:thm:conservation})} & exact & --- \\
Tower compatibility: $\LUV\approx0.0557\,\MPl$, $n_{\rm UV}\approx73.6$ & \textcolor{darkgreen}{Proved (Cor.~\ref{P7:cor:luv_tower})} & exact & --- \\
Cassini Vainshtein radius topologically determined by $Q=2$ & \textcolor{darkgreen}{Proved (Cor.~\ref{P7:cor:vainshtein_topology})} & exact & --- \\
EFT self-consistency below $\LUV$ & \textcolor{darkgreen}{Proved (Prop.~\ref{P7:prop:eft_naturalness})} & exact & --- \\
RG fixed point $\zeta=\vp$ (Derrick scale ratio) & \textcolor{darkgreen}{Derived (Paper~VI~\cite{Paper6})} & exact & --- \\
$\rinfty=\vp/\beta$ (self-consistent) & \textcolor{darkgreen}{Derived} & exact & --- \\
$w=-1$ (frozen field) & \textcolor{darkgreen}{Derived} & exact & --- \\
Acoustic Kerr horizon $=$ Boyer--Lindquist & \textcolor{darkgreen}{Derived} & exact & --- \\
Exact rotating condensate background $\nabla^2\Phi_{\rm Kerr}=0$ & \textcolor{darkgreen}{Proved (Thm~\ref{P7:thm:rotating})} & exact & --- \\
$T_H$ exact GR at leading order & \textcolor{darkgreen}{Derived} & exact & --- \\
Sign of $\delta\rho<0$ near mass & \textcolor{darkgreen}{Derived} & exact & --- \\
WEP at one-loop (geometric) & \textcolor{darkgreen}{Derived} & exact & --- \\
$c_s = 1/\vp$ (exact golden-ratio sound speed) & \textcolor{darkgreen}{Proved (Thm~\ref{P7:thm:cs})} & exact & --- \\
Condensate pressureless on galactic scales ($\lambda_J\gg r_{\rm halo}$) & \textcolor{darkgreen}{Proved (Prop.~\ref{P7:prop:CDM_limit})} & exact & --- \\
$G_{\rm eff}(r)/G_N\leq1+6\times10^{-5}$ inside MW halo & \textcolor{darkgreen}{Proved (Prop.~\ref{P7:prop:Geff})} & exact & --- \\
Adiabatic condensate ICs with $\mathcal{P}_s=2.10\times10^{-9}$ & \textcolor{darkgreen}{Proved (Prop.~\ref{P7:prop:IC_condensate})} & exact & --- \\[4pt]
\SetCell[c=4]{l}{\textit{Tier 1 (continued): Falsifiable predictions awaiting observation}} \\
BAO shift: $k_{\rm BAO}^{\rm cond}=\vp\cdot k_{\rm BAO}^{\rm CDM}\approx0.0346\,\mathrm{Mpc}^{-1}$ & \textcolor{darkgreen}{Proved (Sec.~\ref{P7:sec:transfer})} & $0.0346\,\mathrm{Mpc}^{-1}$ & DESI~DR2, Euclid \\
$r=0.0036$ (tensor-to-scalar) & \textcolor{darkgreen}{Proved (Prop.~\ref{P7:prop:Treh})} & $0.0036$ & CMB-S4 ($r>0.0036$ falsifies) \\
Lorentz violation in voids: $\delta c/c\sim10^{-3}$ & Derived (Sec.~\ref{P7:sec:gravity}) & $\sim10^{-3}$ & VLBI astrometry \\
GW speed $|v_{\rm GW}-c|/c<10^{-33}$ (standard) & Derived & $<10^{-33}$ & $\checkmark$ (GW170817) \\
Hawking WEP instanton $\eta_A=e^{-AS_1}$, testable at STE-QUEST & \textcolor{darkgreen}{Exact (Paper~V~\cite{Paper5})} & $\eta\approx2\times10^{-14}$ & STE-QUEST \\
$w_a=-3(1+w_0)$ (thawing relation) & \textcolor{darkgreen}{Proved (Thm~\ref{P7:thm:thawing})} & falsifiable & $0.80\sigma$ (DESI~DR1) \\
\SetCell[c=4]{l}{\textit{Tier 2: Quantitative successes ($<15\%$)}} \\
$\Lobs/\MPl^4 \approx 10^{-122}$ & Derived & $\sim10^{-122}$ & $1$ order (in exp.) \\
$\rinfty$ consistency & Derived & $10\vp\,\MPl^3$ & $3.4\%$ \\
$n_s = 0.9654$ (Starobinsky, $N_e=57.73$, grav.\ reheating) & \textcolor{darkgreen}{Proved (Prop.~\ref{P7:prop:Treh})} & $0.9654$ & $0.12\sigma$ \\
$r = 0.0036$ (Starobinsky, $N_e=57.73$, grav.\ reheating) & \textcolor{darkgreen}{Proved (Prop.~\ref{P7:prop:Treh})} & $0.0036$ & $\checkmark$ \\
$n_s = 0.9667$ (preliminary, $N_e=60$, no reheating) & Derived & $0.9667$ & $0.4\sigma$ \\
$r = 0.0033$ (preliminary, $N_e=60$, no reheating) & Derived & $0.0033$ & $\checkmark$ \\
$T_H$ correction $+4.2\times10^{-5}$ & Derived & pred. & unmeasured \\
$m_\xi\approx H_0$ & Derived & $\approx H_0$ & $\checkmark$ \\
WEP instanton $\eta_A=e^{-AS_1}$, $S_1\approx0.108$ & \textcolor{darkgreen}{Exact (Paper~V~\cite{Paper5})} & exact for $A\gtrsim277$ & $\checkmark$ \\
Condensate-scale $m_q/m_\ell = e^{n_q\pi/9}$, $n_q\in\mathbb{Z}$ & \textcolor{darkgreen}{Exact (Paper~V~\cite{Paper5})} & exact & QCD running \\
Lepton T-phase uniqueness \& non-integrality ($\pi/\ln\vp\notin\mathbb{Q}$) & \textcolor{darkgreen}{Proved (Thm~5.4, Prop~5.5 of Paper~VI~\cite{Paper6})} & exact & --- \\
$\sinW=3/(8\vp)=0.2318$ & \textcolor{darkgreen}{Unconditional (Cor.~5.7 of~\cite{Paper1})} & $0.2318$ & $0.23\%$ \\
$\alpha^{-1}=137.035\,998$ (unconditional four-term) & \textcolor{darkgreen}{Recalled from~\cite{Paper4}} & $137.035\,998$ & $5\times10^{-7}\%$ \\
$m_e/\Lcond = \vp^{20}e^{4\pi-3/400-\aem/\pi}$ (scheme corrected) & \textcolor{darkgreen}{Paper~IV~\cite{Paper4}, Thm~8.11} & $0.5110\,\mathrm{MeV}$ & $0.013\%$ \\
$m_\mu/m_\tau = e^{-9\pi/10} = e^{-2\pi(Q-1)\Delta T_{23}}$ & \textcolor{darkgreen}{Exact formula; mechanism open (Papers~IV,V~\cite{Paper4,Paper5})} & $0.05917$ & $0.50\%$ \\
$T_3 = \Delta T_{12}(1-1/Q)$,\ equiv.\ $T_1-T_2-T_3 = 1/120$ & \textcolor{darkgreen}{Exact algebraic identity (Papers~III,IV~\cite{Paper3,Paper4})} & exact & exact \\
Single-input framework ($\TCMB$ only) & \textcolor{darkgreen}{Proved (Paper~XII~\cite{Paper12})} & exact chain & $0.013\%$ \\
$\mathcal{P}_s$ see-saw ($N_e=60$, no $R_0$ corr.) & Prop.~\ref{P7:prop:seesaw},\,Thm~\ref{P7:thm:Ps_seesaw} & $2.82\times10^{-9}$ & $\times1.34$ \\
$\mathcal{P}_s$ see-saw ($N_e=60$, with $R_0$ corr.) & Prop.~\ref{P7:prop:seesaw_correction} & $2.27\times10^{-9}$ & $\times1.08$ \\
$\mathcal{P}_s$ see-saw ($N_e=57.73$, grav.\ reheating) & \textcolor{darkgreen}{Proved (Prop.~\ref{P7:prop:Treh})} & $2.10\times10^{-9}$ & exact \\
Cassini via Vainshtein, $r_V\approx9.3\times10^6$\,AU & \textcolor{darkgreen}{Proved (Thm~\ref{P7:thm:cassini}); $|\gamma_{\rm PPN}-1|=2.2\times10^{-5}$} & $\checkmark$ & $\checkmark$ \\
$g_W\approx0.592$ & from~\cite{Paper1} & $0.592$ & $9\%$ \\
$m_H\approx111\,\mathrm{GeV}$ (Superseded) & from~\cite{Paper1} & $111\,\mathrm{GeV}$ & $11\%$ \\[4pt]
$m_H \approx 125.11\,\mathrm{GeV}$ & from~\cite{Paper14} using the framework value $\vEW=246.24\,\mathrm{GeV}$ & $125.11\,\mathrm{GeV}$ &  $0.08\sigma$ \\
$m_H \approx 125.10\,\mathrm{GeV}$ & from~\cite{Paper14} using measured $\vEW=246.22\,\mathrm{GeV}$ & $125.10\,\mathrm{GeV}$ &  $0.13\sigma$ \\
$\lambda_{\mathrm{quartic}} = 0.12907$ & from~\cite{Paper14}: $0.12907$ exact & $0.12907$ & $0.02\%$ \\
\SetCell[c=4]{l}{\textit{Tier 3: Approximate ($15\%$--factor 2)}} \\
$\lambda_{\mathrm{quartic}}\approx0.102$ (Superseded) & from \cite{Paper1}& $0.102$ & $21\%$ \\
$m_\nu \approx \vEW\,e^{-15\pi/2}\cdot g_W^2/(16\pi^2)$ (Prop.~\ref{P7:prop:neutrino}) & \textcolor{darkgreen}{No free parameters; vertex ($\chi_6$, Paper~V) and loop exact} & $36\,\mathrm{meV}$ & $28\%$ \\[4pt]
\SetCell[c=4]{l}{\textit{Tier 4: Known gaps}} \\
$\sinW$ (naive route) $=0.100$ vs.\ $0.231$ & \textcolor{darkred}{Gap: naive identification fails; full derivation in~\cite{Paper1}} & --- & $57\%$ \\
\bottomrule
\end{longtblr}
\end{center}
%══════════════════════════════════════════════════════════════════════════════
\section{Open Problems}
\label{P7:sec:open}
%══════════════════════════════════════════════════════════════════════════════

\begin{enumerate}[noitemsep]

\item \textbf{[Experimental] Dark matter: condensate perturbations and structure formation --- complete.}
  All three computations of Section~\ref{P7:sec:dark_matter} are now proved.
  \textbf{(i)}~$c_s=1/\vp$ exactly (Theorem~\ref{P7:thm:cs}).
  \textbf{(ii)}~Initial conditions are adiabatic with amplitude
  $\mathcal{P}_s = 2.10\times10^{-9}$ (Proposition~\ref{P7:prop:IC_condensate}),
  fixed by the framework with no additional free parameter.
  \textbf{(iii)}~The condensate is pressureless on galactic scales
  ($\lambda_J\approx28\,000\,\mathrm{Mpc}\gg r_{\rm halo}$, Proposition~\ref{P7:prop:CDM_limit})
  and $G_{\rm eff}(r)=G_N$ to $6\times10^{-5}$ at rotation-curve radii
  (Proposition~\ref{P7:prop:Geff}).
  Flat rotation curves follow by the same CDM mechanism.
  The primary new prediction is the BAO peak shift
  $k_{\rm BAO}^{\rm cond}=\vp\cdot k_{\rm BAO}^{\rm CDM}\approx0.0346\,\mathrm{Mpc}^{-1}$,
  falsifiable with DESI~DR2 and Euclid.

\item \textbf{[Experimental] DESI dynamical dark energy.}
  The thawing relation $w_a=-3(1+w_0)$ is proved (Theorem~\ref{P7:thm:thawing})
  and agrees with DESI~DR1 within $0.80\sigma$.
  The open question is observational: whether DESI~DR2 or future surveys
  confirm $w_0\neq-1$ at $>3\sigma$, which would constitute direct evidence
  for the thawing soup scalar over $\Lambda\mathrm{CDM}$.

\item \textbf{[Resolved in Paper~XII~\cite{Paper12}] Prove $\vp^9\sqrt{\Ja\Jfour} = Q$ exactly
(the profile self-consistency conjecture; Papers~II--III~\cite{Paper2,Paper3}).}

  \textbf{This problem does not block any main result of Papers~I--VII.}
  All theorems and predictions in the series hold with the two experimental
  inputs ($\TCMB$ and $m_e$) independently of whether this conjecture is proved.
  Proving it would close the $0.04\%$ thin-torus residual of Paper~IV and
  reduce the framework to a single input ($\TCMB$), but does not alter any
  physical prediction. Resolved in Paper~XII~\cite{Paper12}.

  The problem has three distinct layers that must not be conflated.

  \textbf{Layer 1: $\bz\cdot\rCMB = Q$ --- already proved exactly.}
  The identity $\rCMB/\Lcond^4 = Q = 10$ is algebraically exact
  (proved in Section~\ref{P7:sec:TCMB_geometric}).
  There is nothing further to prove here.  The $0.12\%$ deviation
  in Paper~I arises from $\vp^9\sqrt{\Ja\Jfour}$
  (Papers~II--III~\cite{Paper2,Paper3}), not from $\bz\cdot\rCMB$.

  \textbf{[Resolved in Paper~XII~\cite{Paper12}] Layer 2: $\vp^9\sqrt{\Ja\Jfour} = Q$ --- prior open conjecture~\cite{Paper2,Paper3}.}
  Paper~I's solver gives $\vp^9\sqrt{\Ja\Jfour} = 10.012$ at $h=0.12$.
  The quantity $\vp^9\sqrt{\Ja\Jfour}$ is \emph{scale-invariant}
  (both $\Ja\sim\mathrm{length}$ and $\Jfour\sim\mathrm{length}^{-1}$, so
  their product is dimensionless), depending only on the \emph{shape}
  of the profile, not its overall scale.
  This is a genuine constraint on the EL saddle: it says the condensate
  profile shape exactly encodes its topological charge.
  Paper~I states all five solver gaps are ``consistent in magnitude with
  a single $O(h^2)$ discretisation error'' and provides strong numerical
  evidence the $0.12\%$ closes in the continuum, but does not prove it.
  Papers~I--III prove the related gaps analytically ($J_4/J_{2a}\to
  2^{4/3}/\vp^5$, $V^*\to\vp$), but these fix only the \emph{ratio}
  $J_4/J_{2a}$, not the absolute value $\Ja$.
  The analytical proof requires showing
  $J_{2a,\mathrm{exact}} = Q/(\vp^{13/2}\cdot 2^{2/3}\cdot(2\pi)^2 R_0)$
  at the exact continuum saddle. Corollary~5.10. The analytical proof is closed in Paper~XII~\cite{Paper12}.

  \textbf{[Resolved in Paper~XII~\cite{Paper12} \& Paper~XIV~\cite{Paper14}] Proof route (both gaps closed).}
  The paper \emph{``BPS Structure and Fixed-Point Theorem''}
  (Paper~I~\cite{Paper1}, Section~2.5 therein) derives a linear
  equation for $\Ja$ from the second Pohozaev identity with
  radial multiplier $r\partial_r f$.
  The identity $P_{\Kfb}=2\vp P_{\Ja}$ is proved there as an exact
  algebraic consequence of $\vp^2=\vp+1$ (no profile assumption needed),
  closing the first of the two gaps identified in the original approach.
  The one remaining step is to prove The exact thick-torus Hopf flux
  $\mathcal{B}(R_0=3) = 2Q(1+2^{4/3})/(\vp^{23/2}\cdot 2^{2/3})\approx0.175$
  at the physical condensate radius (Paper~III~\cite{Paper3}, equation~62)
  requires the full 2D profile at the EL saddle, since the thin-torus
  formula $4\pi^2 Q_{\mathrm{math}}/R_0$ is not valid at $R_0=3$
  (thick-torus corrections are $O(1)$).
  The profile normalisation is proved in Paper~XII~\cite{Paper12};
  the exact closed-form thick-torus profile integrals are derived in Paper~XIV~\cite{Paper14}, reducing the $\vEW$ residual from
  $0.69\%$ to $0.01\%$.

  \textbf{Layer 3: Can $\TCMB$ be primary? No --- proved by scale invariance.}
  Proposition~\ref{P7:prop:scale_modulus} (Section~\ref{P7:sec:TCMB_geometric})
  proves that the condensate has a one-parameter family of solutions
  parametrised by $\Lcond$, so $\TCMB$ is not determined by the
  condensate dynamics.  This is a structural result, not a gap.

  \textbf{[Resolved in Paper~XII~\cite{Paper12}] What proving Layer 2 would imply.}
  If $\vp^9\sqrt{\Ja\Jfour}=Q$ is proved analytically~\cite{Paper2,Paper3}:
  \begin{itemize}[noitemsep]
    \item The $0.04\%$ thin-torus spoke residual closes (Paper~IV).
    \item The two routes of Paper~I Corollary~5.10 agree exactly.
    \item Combined with $\rCMB/\Lcond^4=Q$ (exact), this confirms the
      full self-consistency of the condensate at the exact saddle.
  \end{itemize}
  It would reduce the two-input framework ($\TCMB$ and $m_e$) to
  a single experimental input ($\TCMB$ only), but does \emph{not}
  imply $\TCMB$ can be primary (scale invariance prevents this). Proof is closed in Paper~XII~\cite{Paper12}.

  \textbf{[Resolved in Paper~XII~\cite{Paper12}] Proof route.}
  Show that the exact continuum EL saddle forces
  $\Ja = Q/(\vp^{13/2}\cdot 2^{2/3})$ through the normalisation
  condition that equates the two routes of Paper~I, Corollary~5.10. Proof is closed in Paper~XII~\cite{Paper12}.

\end{enumerate}

%══════════════════════════════════════════════════════════════════════════════
\section{Summary}
\label{P7:sec:summary}
%══════════════════════════════════════════════════════════════════════════════

Starting from the four axioms of the density-feedback condensate ---
whose Hopfion vacuum was proved unique in~\cite{Paper1} ---
we derived Newtonian gravity from gradient flux imbalance, the
Schwarzschild and Kerr acoustic metrics with $G_{\mu\nu}=0$ in the
exterior, a one-loop induced Einstein--Hilbert term fixing the
Planck mass via $\MPl^2=\LUV^2/(32\pi^2)$, and the
cosmological constant $\Lobs=\Lambda_0 e^{-\vp^7/\beta}\approx10^{-122}\MPl^4$.
The full dynamical Einstein field equations $G_{\mu\nu}=-\Lambda_{\rm phys}\,g_{\mu\nu}$
are derived from the condensate effective action (Theorem~\ref{P7:thm:einstein}):
the k-essence action reduces to the Einstein--Hilbert action exactly in the frozen
attractor limit, with the effective gravitational strength
$F(\rinfty)=\MPl^2(39\vp+25)/2$ --- an exact algebraic consequence of the golden ratio.
The scalar field equation and energy-momentum conservation $\nabla^\mu T_{\mu\nu}=0$
are proved (Theorem~\ref{P7:thm:conservation}), completing the self-consistent
scalar-tensor structure.
The exact rotating condensate background is derived as the Kerr mass multipole
series $\Phi_{\mathrm{Kerr}}=M/r - Ma^2P_2/(2r^3)+\ldots$, satisfying
$\nabla^2\Phi_{\mathrm{Kerr}}=0$ exactly (Theorem~\ref{P7:thm:rotating}).
The one-loop Planck mass formula implies $\LUV\approx\MPl/(4\pi\sqrt{2})\approx0.0563\,\MPl\approx\MPl/\vp^6\approx0.0557\,\MPl$,
and the corresponding bosonic tower level $n_{\rm UV}\approx73.6$ is non-integer,
confirming that $\LUV$ is set by a continuous gravitational condition, not a discrete
tower step (Corollary~\ref{P7:cor:luv_tower}; connects to Paper~VI~\cite{Paper6}).
The SU(3) colour structure and QCD predictions are treated in
the companion Paper~V~\cite{Paper5}.

The Cassini PPN tension is proved resolved (Theorem~\ref{P7:thm:cassini}):
$m_\xi^2=\vp^{20}\Lambda_{\mathrm{obs}}\MPl^4$ and $\oBD=\vp^{12}/(3\beta)$
combine to give $r_V\approx9.3\times10^6\,\mathrm{AU}$, yielding
$|\gamma_{\mathrm{PPN}}-1|\approx2.2\times10^{-5}$ below the Cassini bound,
with no free parameter.
The topological origin of this Vainshtein radius is also resolved
(Corollary~\ref{P7:cor:vainshtein_topology}): the chain
$Q_{\mathrm{Hopfion}}=2 \to Q_{\mathrm{condensate}}=10 \to m_\xi^2\propto\vp^{20}$
establishes $r_V$ as a topological invariant of the $Q=2$ condensate,
closing the open problem stated in Paper~I~\cite{Paper1}.

Inflationary predictions $n_s=0.9654$ (within $0.12\sigma$ of Planck~2018)
and $r=0.0036$ follow from the Higgs/Starobinsky branch of the
non-minimal coupling $\xNMC R$, at $N_e=57.73$ from gravitational reheating.  The CMB power spectrum normalisation
is proved exactly (Propositions~\ref{P7:prop:seesaw}--\ref{P7:prop:Treh}):
the gravitational see-saw scale $M_{\mathrm{inf}}=\LUV^2/\MPl
\approx3.79\times10^{16}\,\mathrm{GeV}$ gives $\mathcal{P}_s=2.82\times10^{-9}$
at $N_e=60$ (factor $1.34$); the $O(R_0^{-2})$ correction reduces this to $1.08$;
and gravitational reheating ($T_{\mathrm{reh}}\sim10^8\,\mathrm{GeV}$,
no new parameter) gives $N_e=57.73$ and closes the factor to exactly $1.00$,
with $n_s=0.9654$ consistent with Planck~2018 to $0.1\sigma$.
Dark energy satisfies $w=-1$ in
the frozen limit, with the thawing prediction $w_a=-3(1+w_0)$
proved (Theorem~\ref{P7:thm:thawing}) and consistent at $0.80\sigma$ with DESI~DR1.

The condensate scalar field acts as dark matter without introducing a new particle.
The exact golden-ratio sound speed $c_s=1/\vp$ is proved in Theorem~\ref{P7:thm:cs},
following exactly from $\beta\rinfty=\vp$ and $1+\vp=\vp^2$.
The condensate is pressureless to one part in $10^{10}$ on galactic scales
($\lambda_J\approx2\times10^4\,\mathrm{Mpc}\gg r_{\rm halo}$,
Proposition~\ref{P7:prop:CDM_limit}), and the $G_{\rm eff}$ correction inside the
Vainshtein sphere of a Milky-Way galaxy is $\lesssim6\times10^{-5}$
(Proposition~\ref{P7:prop:Geff}).
Rotation curves are therefore flat by the same CDM gravitational-instability
mechanism, with the condensate forming an NFW-like halo.
Initial conditions are adiabatic with amplitude $\mathcal{P}_s=2.10\times10^{-9}$
fixed by the inflationary sector (Proposition~\ref{P7:prop:IC_condensate}),
with no additional free parameter.
The primary new prediction distinguishing the condensate from standard CDM
is the BAO peak shift by factor $\vp$ in the matter power spectrum:
$k_{\rm BAO}^{\rm cond} = \vp\cdot k_{\rm BAO}^{\rm CDM}\approx0.0346\,\mathrm{Mpc}^{-1}$,
falsifiable with DESI~DR2 and Euclid.

The Hawking temperature is recovered exactly at leading order, with
a soup correction $\delta T_H/T_H\approx+4.2\times10^{-5}$ as a
specific model prediction.  The WEP holds exactly at one-loop and is
protected perturbatively by the shift symmetry of the canonical field;
the non-perturbative instanton contribution gives $\eta\approx2\times10^{-14}$,
below MICROSCOPE but testable at STE-QUEST.
The neutrino mass is derived via the WZW radiative mechanism
(Proposition~\ref{P7:prop:neutrino}): the $\chi_6$ Node~6 primary
of the $2I$ orbifold spectrum ($h=3/8$, Paper~V Theorem~11.7)
provides the $\Delta L=2$ vertex and seed
$m_\mathrm{seed}=\vEW\,e^{-15\pi/2}\approx14.4\,\mathrm{eV}$;
one $W$-boson loop with factor $g_W^2/(16\pi^2)=2\alpha\vp/(3\pi)$
gives $m_\nu\approx36\,\mathrm{meV}$, within $28\%$ of the atmospheric
scale (Tier~3, no free parameters).
The loop is regularised by the condensate's physical UV cutoff $\LUV$
via the same Seeley--DeWitt method that fixes $\MPl$; corrections
are $O((M_W/\LUV)^2)\approx10^{-32}$.

The framework connects the Hopfion vacuum of Papers~I--VI to a
self-consistent gravitational and cosmological sector, with all
significant gaps honestly accounted for in Section~\ref{P7:sec:tier}.
The EFT below $\LUV$ is proved self-consistent across all four physical
sectors --- inflation, dark energy, gravity, and the Cassini/WEP
sectors --- with every loop correction suppressed by $(E/\LUV)^2$
(Proposition~\ref{P7:prop:eft_naturalness}).

Paper~IV~\cite{Paper4} (Theorem~4.7, Verlinde cancellation) establishes
that the ratio $m_e/\TCMB$ is a dimensionless geometric invariant of the
$Q=2$ icosahedral condensate, equal to
$(\pi^2/150)^{1/4}\cdot\vp^{20}\cdot e^{(1600\pi-3)/400}$
with no free parameters and $0.22\%$ accuracy ($O(R_0^{-2})$ thick-torus correction).
The framework therefore operates with one experimental inputs,
$\TCMB$, and zero free parameters.
Proving the $O(R_0^{-2})$ residual exactly --- equivalent to
$\vp^9\sqrt{\Ja\Jfour}=Q$ (Papers~II--III~\cite{Paper2,Paper3}) --- reduces the input count to one, and is resolved in Paper~XII~\cite{Paper12}.
The equation $\bst\cdot\rCMB = Q$ ($0.12\%$ from exact in the profile
integrals) reflects the approximate self-consistency of the numerical
saddle; the algebraic identity $\rCMB/\Lcond^4 = Q$ holds exactly.
A scale-invariance argument (Section~\ref{P7:sec:open}) shows that $\TCMB$
is not derivable from the condensate dynamics: the condensate has a
one-parameter family of solutions at every scale, and $\TCMB$ fixes
the physical one.
Three observational questions remain open: whether DESI~DR2 confirms
$w_0\neq-1$ at $>3\sigma$; whether DESI~DR2 or Euclid detect the BAO peak
shift to $k_{\rm BAO}^{\rm cond}\approx0.0346\,\mathrm{Mpc}^{-1}$;
and whether $\vp^9\sqrt{\Ja\Jfour}=Q$ exactly
in the continuum~\cite{Paper2,Paper3}.
The scale modulus $\Lcond$ remains unfixed by the condensate dynamics;
$\TCMB$ is the single experimental input of the framework.

%─── References ───────────────────────────────────────────────────────────────
\begin{thebibliography}{99}

\bibitem{Paper1}
F.~Manfredi,
\textit{The Density-Feedback Faddeev--Niemi Hopfion:
  Exact Algebraic Predictions for Standard-Model Parameters},
Zenodo (2026).
\href{https://doi.org/10.5281/zenodo.19342027}{doi:10.5281/zenodo.19342027}

\bibitem{Paper2}
F.~Manfredi,
\textit{BPS Structure and Fixed-Point Theorem for the
  Density-Feedback Faddeev--Niemi Hopfion},
Zenodo (2026).
\href{https://doi.org/10.5281/zenodo.19363491}{doi:10.5281/zenodo.19363491}

\bibitem{Paper3}
F.~Manfredi,
\textit{Two Constants from One Knot: The Fine Structure Constant and
  the Linking Scale as Exact Predictions of the Icosahedral WZW Condensate},
Zenodo (2026).
\href{https://doi.org/10.5281/zenodo.19478629}{doi:10.5281/zenodo.19478629}

\bibitem{Paper4}
F.~Manfredi,
\textit{The Hopf Spoke, WZW Fermion Mass Renormalization,
  and Three- and Four-Term Formulas for the Fine Structure Constant},
Zenodo (2026).
\href{https://doi.org/10.5281/zenodo.19504729}{doi:10.5281/zenodo.19504729}

\bibitem{Paper5}
F.~Manfredi,
\textit{Colour Confinement, the QCD Scale, and Hadronic Structure
  from the Density-Feedback Hopfion},
Zenodo (2026).
\href{https://doi.org/10.5281/zenodo.19638857}{doi:10.5281/zenodo.19638857}

\bibitem{Paper6}
F.~Manfredi,
\textit{The Golden-Ratio Tower: Fermion Scale Hierarchy from the Hopfion RG Fixed Point},
Zenodo (2026).
\href{https://doi.org/10.5281/zenodo.19646945}{doi:10.5281/zenodo.19646945}

\bibitem{Paper12}
F.~Manfredi,
\textit{The Profile Normalisation Conjecture ---
  Proof via CMB Energy Identity and Two-Route Equivalence},
Zenodo (2026).
\href{https://doi.org/10.5281/zenodo.20210460}{doi:10.5281/zenodo.20210460}

\bibitem{Paper14}
F.~Manfredi,
\textit{The Thick-Torus Profile Correction to the Density-Feedback Hopfion},
Zenodo (2026).
\href{https://doi.org/10.5281/zenodo.20479790}{doi:10.5281/zenodo.20479790}

\bibitem{Unruh1981}
W.~G. Unruh,
\textit{Experimental black-hole evaporation?},
Phys.\ Rev.\ Lett.\ \textbf{46}, 1351 (1981).
\href{https://doi.org/10.1103/PhysRevLett.46.1351}{doi:10.1103/PhysRevLett.46.1351}

\bibitem{Cassini2003}
B.~Bertotti, L.~Iess, and P.~Tortora,
\textit{A test of general relativity using radio links with the
Cassini spacecraft},
Nature \textbf{425}, 374 (2003).
\href{https://doi.org/10.1038/nature01997}{doi:10.1038/nature01997}

\bibitem{Vainshtein1972}
A.~I. Vainshtein,
\textit{To the problem of nonvanishing gravitation mass},
Phys.\ Lett.\ B \textbf{39}, 393 (1972).
\href{https://doi.org/10.1016/0370-2693(72)90147-5}{doi:10.1016/0370-2693(72)90147-5}

\bibitem{Dvali2000}
G.~Dvali, G.~Gabadadze, and M.~Porrati,
\textit{4D gravity on a brane in 5D Minkowski space},
Phys.\ Lett.\ B \textbf{485}, 208 (2000).
\href{https://doi.org/10.1016/S0370-2693(00)00669-9}{doi:10.1016/S0370-2693(00)00669-9}

\bibitem{BabichevDeffayet2013}
E.~Babichev and C.~Deffayet,
\textit{An introduction to the Vainshtein mechanism},
Class.\ Quantum Grav.\ \textbf{30}, 184001 (2013).
\href{https://doi.org/10.1088/0264-9381/30/18/184001}{doi:10.1088/0264-9381/30/18/184001}

\bibitem{Jacobson1995}
T.~Jacobson,
\textit{Thermodynamics of spacetime: The Einstein equation of state},
Phys.\ Rev.\ Lett.\ \textbf{75}, 1260 (1995).
\href{https://doi.org/10.1103/PhysRevLett.75.1260}{doi:10.1103/PhysRevLett.75.1260}

\bibitem{LIGO_GW170817}
B.~P. Abbott et al.\ (LIGO Scientific and Virgo Collaborations),
\textit{GW170817: Observation of gravitational waves from a binary
neutron star inspiral},
Phys.\ Rev.\ Lett.\ \textbf{119}, 161101 (2017).
\href{https://doi.org/10.1103/PhysRevLett.119.161101}{doi:10.1103/PhysRevLett.119.161101}

\bibitem{Planck2018}
N.~Aghanim et al.\ (Planck Collaboration),
\textit{Planck 2018 results.\ VI.\ Cosmological parameters},
Astron.\ Astrophys.\ \textbf{641}, A6 (2020).
\href{https://doi.org/10.1051/0004-6361/201833910}{doi:10.1051/0004-6361/201833910}

\bibitem{DESI2024}
A.~G. Adame et al.\ (DESI Collaboration),
\textit{DESI 2024 VI: Cosmological constraints from the measurements
of baryon acoustic oscillations},
\href{https://doi.org/10.1088/1475-7516/2025/02/021}{doi:10.1088/1475-7516/2025/02/021}

\bibitem{MICROSCOPE2022}
P.~Touboul et al.\ (MICROSCOPE Collaboration),
\textit{MICROSCOPE mission: Final results of the test of the equivalence
principle},
Phys.\ Rev.\ Lett.\ \textbf{129}, 121102 (2022).
\href{https://doi.org/10.1103/PhysRevLett.129.121102}{doi:10.1103/PhysRevLett.129.121102}

\bibitem{KhouryWeltman2004}
J.~Khoury and A.~Weltman,
\textit{Chameleon fields: Awaiting surprises for tests of gravity in space},
Phys.\ Rev.\ Lett.\ \textbf{93}, 171104 (2004).
\href{https://doi.org/10.1103/PhysRevLett.93.171104}{doi:10.1103/PhysRevLett.93.171104}

\bibitem{PDG2024}
R.~L. Workman et al.\ (Particle Data Group),
\textit{Review of Particle Physics},
Prog.\ Theor.\ Exp.\ Phys.\ \textbf{2022}, 083C01 (2022);
updated 2024.
\href{https://doi.org/10.1093/ptep/ptac097}{doi:10.1093/ptep/ptac097}

\bibitem{Koide1983}
Y.~Koide,
\textit{A fermion-boson composite model of quarks and leptons},
Phys.\ Lett.\ B \textbf{120}, 161--165 (1983).
\href{https://doi.org/10.1016/0370-2693(83)90644-5}{doi:10.1016/0370-2693(83)90644-5}

\bibitem{McKay1980}
J.~McKay,
\textit{Graphs, singularities, and finite groups},
Proc.\ Symp.\ Pure Math.\ \textbf{37}, 183--186 (1980).
\href{https://doi.org/10.1090/pspum/037}{doi:10.1090/pspum/037}

\end{thebibliography}

\end{document}